\documentclass[12pt, a4paper]{report}
\usepackage{graphicx}
\usepackage{float}
\usepackage[hidelinks]{hyperref}
\usepackage{array}
\usepackage{longtable}
\usepackage{booktabs}
\usepackage{polyglossia}
\usepackage{geometry}
\geometry{a4paper, margin=1in}

\setmainlanguage[numerals=maghrib]{arabic}
\setotherlanguage{english}
\newfontfamily\arabicfont[Script=Arabic]{Arial} % You can change this to Amiri or Noto Sans Arabic on Overleaf

\begin{document}
\vspace{0.5cm}\hrule\vspace{0.5cm}
\begin{center}

\textbf{وزارة التعليم العالي} \\

\textbf{أكاديمية الدلتا للعلوم والتكنولوجيا} \\

\textbf{معهد الدلتا العالي لنظم المعلومات الإدارية والمحاسبية بالمنصورة} \\

\vspace{0.5cm}\hrule\vspace{0.5cm}

\begin{figure}[H]
\centering
\includegraphics[width=0.8\textwidth]{extracted_images/img_00.png}
\caption{شعار المعهد}
\end{figure}

\vspace{0.5cm}\hrule\vspace{0.5cm}

\chapter*{مشروع رقم (44)}
\addcontentsline{toc}{chapter}{مشروع رقم (44)}
\section*{منصة الكترونية لدعم المؤسسات استعداداً لتطبيق معايير الجودة}
\addcontentsline{toc}{section}{منصة الكترونية لدعم المؤسسات استعداداً لتطبيق معايير الجودة}
\subsection*{(Ayn Platform)}
\addcontentsline{toc}{subsection}{(Ayn Platform)}

\vspace{0.5cm}\hrule\vspace{0.5cm}

\textbf{مشروع التخرج مقدم كمتطلب جزئي للحصول على درجة البكالوريوس} \\
\textbf{في تخصص: نظم معلومات الأعمال} \\

\vspace{0.5cm}\hrule\vspace{0.5cm}

\textbf{تحت إشراف:} \\

\textbf{د. أسماء عبد الباسط} \\

\textbf{مساعد الإشراف:} \\

\textbf{م. نورهان مراد} \\

\vspace{0.5cm}\hrule\vspace{0.5cm}

\textbf{2025 – 2026} \\

\end{center}

\vspace{0.5cm}\hrule\vspace{0.5cm}
\newpage

\begin{center}

بسم الله الرحمن الرحيم \\

\vspace{0.5cm}\hrule\vspace{0.5cm}

\section*{﴿ وَقُل رَّبِّ زِدۡنِي عِلۡمٗا ﴾}
\addcontentsline{toc}{section}{﴿ وَقُل رَّبِّ زِدۡنِي عِلۡمٗا ﴾}

\textbf{صدق الله العظيم} \\

\textit{(سورة طه: الآية 114)} \\

\end{center}

\vspace{0.5cm}\hrule\vspace{0.5cm}
\newpage

\begin{center}

\chapter*{إهداء}
\addcontentsline{toc}{chapter}{إهداء}

\end{center}

إلى من أناروا دروبنا بعلمهم ووجهونا بحكمتهم... \\
إلى الذين آمنوا بنا قبل أن نؤمن بأنفسنا... \\

إلى \textbf{أمهاتنا وآبائنا} الذين سهروا الليالي لتكون ليالينا هادئة، ودفعوا بنا نحو النجاح بصبر لا حدود له وعطاء لا ينضب. كل حرف من هذا البحث هو ثمرة من ثمار تضحياتهم. \\

إلى \textbf{أساتذتنا الكرام} الذين لم يبخلوا علينا بعلم ولا بتوجيه، وفتحوا أمامنا آفاقاً لم نكن لنصل إليها بمفردنا. \\

إلى \textbf{أصدقائنا وزملائنا} الذين شاركونا رحلة الدراسة بكل ما فيها من متعة وتعب، ونجاح وتحدٍّ. \\

إلى \textbf{كل من آمن بأن التكنولوجيا يمكن أن تُحدث فارقاً حقيقياً} في منظومة التعليم وجودته. \\

إلى كل هؤلاء نهدي هذا الجهد المتواضع، راجين من الله العلي القدير أن يكون خالصاً لوجهه الكريم، وأن ينفع به. \\

\begin{center}

\textit{"من لا يشكر الناس لا يشكر الله"} \\

\textbf{فريق المشروع} \\

\end{center}

\vspace{0.5cm}\hrule\vspace{0.5cm}
\newpage

\begin{center}

\chapter*{شكر وتقدير}
\addcontentsline{toc}{chapter}{شكر وتقدير}

\textit{الشكر في عظيم الصنع قليل يفري رواء الهمم نهراً بالخيرات يسيل إلى من أعطوا وأجزوا بعطائهم.....} \\

\vspace{0.5cm}\hrule\vspace{0.5cm}

\section*{د/ محمد ربيع ناصر}
\addcontentsline{toc}{section}{د/ محمد ربيع ناصر}
\subsection*{رئيس مجلس إدارة أكاديمية الدلتا للعلوم والتكنولوجيا}
\addcontentsline{toc}{subsection}{رئيس مجلس إدارة أكاديمية الدلتا للعلوم والتكنولوجيا}

\textit{[صورة د. محمد ربيع ناصر]} \\

\end{center}

رائد من رواد التعليم العالي الخاص والمتميز، ونفتخر بأننا حصاد لتعليم متميز قام سيادته بإرساء مبادئه، مما أتاح لنا الفرصة للتعلم والتزود بباكورة مؤسساته بمعهد الدلتا العالي لنظم المعلومات الإدارية والمحاسبية، حيث شملنا بكل الرعاية وقدم لنا الدعم لاكتساب المهارات العلمية والمعرفة الحقيقية التي كانت خير عون لنا. فشكراً له جزيل الشكر. \\

\vspace{0.5cm}\hrule\vspace{0.5cm}
\newpage

\begin{center}

\section*{أ.د/ أحمد أبو الفتوح}
\addcontentsline{toc}{section}{أ.د/ أحمد أبو الفتوح}
\subsection*{عميد معهد الدلتا العالي لنظم المعلومات الإدارية والمحاسبية}
\addcontentsline{toc}{subsection}{عميد معهد الدلتا العالي لنظم المعلومات الإدارية والمحاسبية}

\textit{[صورة أ.د. أحمد أبو الفتوح]} \\

\end{center}

نتقدم بأسمى آيات الشكر إلى أ.د/ أحمد أبو الفتوح، الذي نفتخر بأن تخرجنا كان تحت عمادته، حيث أعاد لنا ولمعهد الدلتا الريادة من خلال دعمه المتواصل لنا وتطويره لكل ما فيه صالحنا في مخرجات التعليم والتدريب ومتابعته الدؤوبة لنا في كل خطوات التعلم. لن توفيكم كلمات الشكر، فستظل دوما في فكرنا داعين الله لسيادته بدوام التوفيق والسداد والتقدم. \\

\vspace{0.5cm}\hrule\vspace{0.5cm}
\newpage

\begin{center}

\section*{د/ أسماء عبد الباسط}
\addcontentsline{toc}{section}{د/ أسماء عبد الباسط}
\subsection*{المشرف على المشروع}
\addcontentsline{toc}{subsection}{المشرف على المشروع}

\textit{[صورة د. أسماء عبد الباسط]} \\

\end{center}

نتقدم بجزيل الشكر والتقدير والامتنان إلى د/ أسماء عبد الباسط التي تفضلت بإشرافها على هذا المشروع، ولكل ما قدمته لنا من دعم وتوجيه وإرشاد لإتمام هذا العمل على أكمل وجه. لقد كان لملاحظاتها الدقيقة وتوجيهاتها السديدة الأثر الكبير في إخراج هذا المشروع بصورته الراهنة. فلها أسمى عبارات الثناء والتقدير، وندعو الله أن يجزيها عنا خير الجزاء. \\

\vspace{0.5cm}\hrule\vspace{0.5cm}
\newpage

\begin{center}

\section*{م/ نورهان مراد}
\addcontentsline{toc}{section}{م/ نورهان مراد}
\subsection*{مساعد الإشراف}
\addcontentsline{toc}{subsection}{مساعد الإشراف}

\textit{[صورة م. نورهان مراد]} \\

\end{center}

تحية شكر وتقدير منا إلى م/ نورهان مراد على الجهد المبذول والوقت الثمين الذي منحتنا إياه في الإرشاد والمتابعة. مهما كتبنا من عبارات الشكر لا نجد خيراً من الدعاء بأن يوفقها الله ويجازيها عنا خيراً، فكان لتواجدها وتوجيهاتها الأثر الطيب في مسيرة العمل على هذا المشروع. \\

\vspace{0.5cm}\hrule\vspace{0.5cm}
\newpage

\begin{center}

\chapter*{فريق العمل (Team Work)}
\addcontentsline{toc}{chapter}{فريق العمل (Team Work)}

\end{center}

\begin{longtable}{|c|c|c|}
\hline
الصورة & الصورة & الصورة \\ \hline
:---: & :---: & :---: \\ \hline
\textit{[صورة]} & \textit{[صورة]} & \textit{[صورة]} \\ \hline
\textbf{الاسم} & \textbf{الاسم} & \textbf{الاسم} \\ \hline
\textit{[صورة]} & \textit{[صورة]} & \textit{[صورة]} \\ \hline
\textbf{الاسم} & \textbf{الاسم} & \textbf{الاسم} \\ \hline
\textit{[صورة]} & \textit{[صورة]} & \textit{[صورة]} \\ \hline
\textbf{الاسم} & \textbf{الاسم} & \textbf{الاسم} \\ \hline
\end{longtable}

\vspace{0.5cm}\hrule\vspace{0.5cm}
\newpage

\begin{center}

\chapter*{ملخص المشروع (Abstract)}
\addcontentsline{toc}{chapter}{ملخص المشروع (Abstract)}

\end{center}

يقدم هذا المشروع "منصة عين" (Ayn Platform)، وهي منصة إلكترونية سحابية مبتكرة صُممت خصيصاً لدعم المؤسسات التعليمية في استعدادها لتطبيق معايير الجودة والحصول على الاعتمادات الأكاديمية الوطنية والدولية. تهدف المنصة إلى تحويل عملية ضمان الجودة الأكاديمية من عملية يدوية مجهدة ومستهلكة للوقت إلى عملية رقمية ذكية ومؤتمتة، تستفيد من أحدث تقنيات الذكاء الاصطناعي التوليدي (Generative AI) وتقنية RAG (Retrieval-Augmented Generation). \\

تتضمن المنصة عدة وحدات متكاملة أبرزها: مستودع الأدلة الذكي (Evidence Vault) لتخزين وتحليل الوثائق المؤسسية تلقائياً، ومحرك الذكاء الاصطناعي "Horus AI" المدعوم بنموذج Google Gemini 1.5 Pro للإجابة على الاستفسارات الأكاديمية وتحليل الامتثال، ووحدة تحليل الفجوات (Gap Analysis) لتحديد مواطن القصور واقتراح خطط التحسين، إضافة إلى لوحات قيادة (Dashboards) تفاعلية وتحليلات إحصائية لحظية. \\

تم بناء المنصة باستخدام تقنيات متقدمة تشمل: Next.js و React للواجهة الأمامية، و FastAPI و Python للواجهة الخلفية، وقاعدة بيانات PostgreSQL مع إضافة pgvector للبحث الدلالي، وخدمة Supabase للتخزين السحابي، وUpstash Redis للتخزين المؤقت. يتبع النظام معمارية Domain-Driven Design ونظام RBAC للتحكم في الصلاحيات. \\

أثبتت نتائج الاختبار أن المنصة قادرة على تحليل الوثائق وتوليد تقارير الفجوات بدقة عالية، مع توفير واجهة مستخدم سهلة وجذابة تدعم اللغتين العربية والإنجليزية وتعمل على كافة الأجهزة. \\

\textbf{الكلمات المفتاحية:} ضمان الجودة الأكاديمية، الاعتماد المؤسسي، الذكاء الاصطناعي، تحليل الفجوات، RAG، منصة سحابية، NAQAAE، ISO 21001. \\

\vspace{0.5cm}\hrule\vspace{0.5cm}

\begin{english}

\textbf{Abstract:} This project presents "Ayn Platform," an innovative cloud-based electronic platform specifically designed to support educational institutions in their preparation for quality standards compliance and academic accreditation. The platform aims to transform academic quality assurance from a tedious, time-consuming manual process into a smart, automated digital operation leveraging cutting-edge Generative AI and RAG technology. Built with Next.js, FastAPI, PostgreSQL+pgvector, and Google Gemini 1.5 Pro, it offers evidence management, gap analysis, AI-powered assistance, and real-time analytics dashboards. \\

\textbf{Keywords:} Academic Quality Assurance, Institutional Accreditation, Artificial Intelligence, Gap Analysis, RAG, Cloud Platform, NAQAAE, ISO 21001. \\

\end{center}

\newpage
\chapter*{فهرس المحتويات}

\begin{longtable}{|c|c|}
\hline
الموضوع & الصفحة \\ \hline
:--- & ---: \\ \hline
الآية القرآنية & 2 \\ \hline
الإهداء & 3 \\ \hline
شكر وتقدير & 4 \\ \hline
فريق العمل & 8 \\ \hline
ملخص المشروع & 9 \\ \hline
فهرس المحتويات & 11 \\ \hline
قائمة الأشكال & 13 \\ \hline
قائمة الجداول & 14 \\ \hline
[\textbf{الفصل الأول: الإطار العام للمشروع}](Chapter_1.md) & \textbf{15} \\ \hline
[1.1 مقدمة](Chapter_1.md#11-مقدمة) & 16 \\ \hline
[1.2 لمحة عامة ووصف النظام](Chapter_1.md#12-لمحة-عامة-ووصف-النظام) & 18 \\ \hline
[1.3 تعريف المشكلة](Chapter_1.md#13-تعريف-المشكلة) & 21 \\ \hline
[1.4 الأهداف](Chapter_1.md#14-الأهداف) & 23 \\ \hline
[1.5 الدراسات السابقة](Chapter_1.md#15-الدراسات-السابقة) & 25 \\ \hline
[1.6 خطة الدراسة](Chapter_1.md#16-خطة-الدراسة) & 28 \\ \hline
[\textbf{الفصل الثاني: فكرة المشروع}](Chapter_2.md) & \textbf{31} \\ \hline
[2.1 أهداف المشروع](Chapter_2.md#21-أهداف-المشروع) & 32 \\ \hline
[2.2 نطاق وقيود المشروع](Chapter_2.md#22-نطاق-وقيود-المشروع) & 36 \\ \hline
[2.3 أهمية المشروع](Chapter_2.md#23-أهمية-المشروع) & 39 \\ \hline
[2.4 المستخدمون المستهدفون](Chapter_2.md#24-المستخدمون-المستهدفون-والمؤسسات) & 41 \\ \hline
[2.5 القيمة الأساسية المقترحة](Chapter_2.md#25-القيمة-الأساسية-المقترحة) & 44 \\ \hline
[\textbf{الفصل الثالث: الأدوات والتقنيات}](Chapter_3.md) & \textbf{47} \\ \hline
[3.1 تقنيات الواجهة الأمامية](Chapter_3.md#31-التقنيات-المستخدمة-في-الواجهة-الأمامية) & 48 \\ \hline
[3.2 تقنيات الواجهة الخلفية](Chapter_3.md#32-التقنيات-المستخدمة-في-الواجهة-الخلفية) & 51 \\ \hline
[3.3 قواعد البيانات والتخزين المؤقت](Chapter_3.md#33-قواعد-البيانات-والتخزين-المؤقت) & 54 \\ \hline
[3.4 دمج الذكاء الاصطناعي](Chapter_3.md#34-دمج-الذكاء-الاصطناعي) & 57 \\ \hline
[3.5 النشر والتكامل المستمر](Chapter_3.md#35-النشر-والتكامل-المستمر) & 59 \\ \hline
[\textbf{الفصل الرابع: تحليل النظام والهندسة المعمارية}](Chapter_4.md) & \textbf{61} \\ \hline
[4.1 وحدات المنصة وهيكليتها](Chapter_4.md#41-وحدات-المنصة-وهيكليتها) & 62 \\ \hline
[4.2 تدفق البيانات](Chapter_4.md#42-تدفق-البيانات) & 66 \\ \hline
[4.3 تصميم قاعدة البيانات](Chapter_4.md#43-تصميم-قاعدة-البيانات) & 70 \\ \hline
[\textbf{الفصل الخامس: التنفيذ}](Chapter_5.md) & \textbf{75} \\ \hline
[5.1 تطوير الواجهة الأمامية](Chapter_5.md#51-عملية-تطوير-الواجهة-الأمامية) & 76 \\ \hline
[5.2 تطوير الواجهة الخلفية](Chapter_5.md#52-عملية-تطوير-الواجهة-الخلفية) & 79 \\ \hline
[5.3 تنفيذ الميزات الأساسية](Chapter_5.md#53-تنفيذ-الميزات-الأساسية) & 82 \\ \hline
[5.4 تقنيات تحسين الأداء](Chapter_5.md#54-تقنيات-تحسين-الأداء) & 87 \\ \hline
[\textbf{الفصل السادس: تصميم UI/UX}](Chapter_6.md) & \textbf{89} \\ \hline
[6.1 هيكل الواجهة الأساسي](Chapter_6.md#61-هيكل-الواجهة-الأساسي) & 90 \\ \hline
[6.2 رحلة المستخدم](Chapter_6.md#62-رحلة-المستخدم) & 93 \\ \hline
[6.3 تخطيطات لوحة القيادة](Chapter_6.md#63-تخطيطات-لوحة-القيادة-والتحليلات) & 95 \\ \hline
[6.4 قرارات التصميم](Chapter_6.md#64-قرارات-التصميم-الجمالية-والفلسفية) & 97 \\ \hline
[\textbf{الفصل السابع: الاختبار والنتائج}](Chapter_7.md) & \textbf{100} \\ \hline
[\textbf{الفصل الثامن: التطويرات المستقبلية}](Chapter_8.md) & \textbf{110} \\ \hline
[\textbf{الخاتمة والتوصيات}](Chapter_9_Conclusion.md) & \textbf{115} \\ \hline
[\textbf{المراجع والمصادر}](References.md) & \textbf{117} \\ \hline
[\textbf{الملاحق}](Appendix.md) & \textbf{120} \\ \hline
\end{longtable}

\vspace{0.5cm}\hrule\vspace{0.5cm}

\chapter*{قائمة الأشكال}

\begin{longtable}{|c|c|c|}
\hline
رقم الشكل & عنوان الشكل & الصفحة \\ \hline
:---: & :--- & ---: \\ \hline
1-1 & الشعار الرسمي للمنصة & 10 \\ \hline
1-2 & المفهوم الأساسي للنظام & 19 \\ \hline
3-1 & مخطط تدفق التكنولوجيا & 49 \\ \hline
3-2 & معمارية تطبيق RAG & 53 \\ \hline
3-3 & خريطة تكامل الذكاء الاصطناعي & 58 \\ \hline
4-1 & المعمارية العامة للنظام & 62 \\ \hline
4-2 & مخطط تدفق البيانات (DFD) & 67 \\ \hline
4-3 & مخطط علاقة الكيانات (ERD) & 71 \\ \hline
5-1 & واجهة مستودع الأدلة & 77 \\ \hline
5-2 & واجهة مساعد Horus AI & 85 \\ \hline
6-1 & هيكل الواجهة والشريط الجانبي & 91 \\ \hline
6-2 & لوحة القيادة الرئيسية & 96 \\ \hline
\end{longtable}

\vspace{0.5cm}\hrule\vspace{0.5cm}

\chapter*{قائمة الجداول}

\begin{longtable}{|c|c|c|}
\hline
رقم الجدول & عنوان الجدول & الصفحة \\ \hline
:---: & :--- & ---: \\ \hline
1-1 & مقارنة الأنظمة السابقة مع منصة Ayn & 26 \\ \hline
2-1 & أدوار المستخدمين وصلاحياتهم & 42 \\ \hline
2-2 & جدول القيمة الأساسية المقترحة & 44 \\ \hline
3-1 & حزمة التكنولوجيا الكاملة & 48 \\ \hline
4-1 & وحدات المنصة ومسؤولياتها & 63 \\ \hline
7-1 & نتائج الاختبار الوظيفي & 104 \\ \hline
\end{longtable}

\newpage
\chapter*{الفصل الأول: الإطار العام للمشروع (Introduction)}
\addcontentsline{toc}{chapter}{الفصل الأول: الإطار العام للمشروع (Introduction)}

\section*{1.1 مقدمة:}
\addcontentsline{toc}{section}{1.1 مقدمة:}

في عصر تتسارع فيه وتيرة التطور التكنولوجي والتحول الرقمي في شتى المجالات، أصبحت جودة التعليم والاعتماد الأكاديمي ركيزة أساسية تعتمد عليها المؤسسات التعليمية لضمان مكانتها التنافسية وتقديم خدمات تعليمية ترقى إلى المعايير العالمية. إن تحقيق معايير الجودة ليس مجرد استيفاء لمتطلبات إدارية أو ورقية، بل هو عملية مستمرة تتطلب تخطيطاً استراتيجياً، وتنفيذاً دقيقاً، ومتابعة حثيثة لضمان التحسين المستمر في الأداء المؤسسي والأكاديمي. ومع تزايد تعقيد أطر الاعتماد وتعدد متطلباتها، أصبح الاعتماد على الأساليب التقليدية اليدوية في إدارة ملفات الجودة يشكل عبئاً ثقيلاً على كاهل المؤسسات التعليمية، مما يؤدي إلى استنزاف الموارد، وضياع الوقت، وارتفاع احتمالية حدوث الأخطاء البشرية. \\

من هذا المنطلق، تبرز الحاجة الملحة لابتكار حلول تقنية متقدمة تساهم في أتمتة وتسهيل عمليات ضمان الجودة. وهنا يأتي دور "منصة عين" (Ayn Platform)، وهي منصة إلكترونية مبتكرة تم تطويرها خصيصاً لدعم المؤسسات التعليمية في استعدادها لتطبيق معايير الجودة والاعتماد الأكاديمي. تمثل المنصة نقلة نوعية في كيفية إدارة الجودة، حيث تدمج بين أحدث تقنيات الويب والذكاء الاصطناعي التوليدي لتقديم بيئة عمل رقمية متكاملة. \\

\begin{figure}[H]
\centering
\includegraphics[width=0.8\textwidth]{extracted_images/img_01.png}
\caption{شعار المنصة والواجهة الرئيسية المبدئية}
\end{figure}

تعمل منصة Ayn كنظام إدارة شامل يبسط العمليات المعقدة المرتبطة بضمان الجودة الأكاديمية والامتثال. فهي لا تقتصر على كونها مستودعاً رقمياً للملفات، بل تتعدى ذلك لتكون مساعداً ذكياً يحلل الفجوات، ويوجه المؤسسات نحو الخطوات التصحيحية اللازمة. من خلال الاستفادة من قوة محرك الذكاء الاصطناعي "Horus AI"، تتيح المنصة للمؤسسات التعليمية تقييم وضعها الحالي بدقة، والتنبؤ بمكامن الضعف، وبناء خطط تحسين مدعومة بالبيانات. إن هذا المشروع يهدف في جوهره إلى تحويل رحلة الاعتماد الأكاديمي من رحلة شاقة ومعقدة إلى مسار واضح، منظم، وفعال، مما ينعكس إيجاباً على مخرجات العملية التعليمية ككل ويعزز من ثقة المجتمع في المؤسسة التعليمية. \\

\section*{1.2 لمحة عامة ووصف النظام:}
\addcontentsline{toc}{section}{1.2 لمحة عامة ووصف النظام:}

نظام "Ayn" هو عبارة عن منصة SaaS (البرمجيات كخدمة) سحابية، مصممة بأحدث الهيكليات البرمجية لتلبي احتياجات الجامعات، المعاهد، والمدارس التي تسعى للحصول على الاعتمادات الوطنية (مثل NAQAAE في مصر أو NCAAA في السعودية) أو الدولية (مثل ISO 21001). يتكون النظام من مجموعة من الوحدات المترابطة التي تعمل بتناغم لتوفير تجربة مستخدم سلسة وفعالة. \\

\textbf{الوحدات الأساسية للنظام تتضمن:} \\
1. \textbf{مستودع الأدلة الذكي (Evidence Vault):} وهو قلب النظام حيث يتم رفع وتخزين كافة الوثائق والسياسات والأدلة الداعمة. لا يقتصر دور المستودع على التخزين السحابي الآمن فحسب، بل يقوم تلقائياً بتهيئة الملفات (سواء كانت نصوصاً، صوراً، أو ملفات PDF) لتكون قابلة للبحث الدلالي والتحليل عبر تقنيات Vector Embeddings. \\
2. \textbf{محرك الذكاء الاصطناعي (Horus AI):} يمثل العقل المدبر للمنصة. يعتمد على نماذج لغوية متقدمة (LLMs) مثل Gemini 1.5 Pro، ويقوم بمهام متعددة تشمل: القراءة العميقة للأدلة، تصنيفها، مطابقتها مع معايير الجودة المستهدفة، وتقديم استشارات تفاعلية للمستخدمين من خلال واجهة محادثة ذكية. \\
3. \textbf{وحدة تحليل الفجوات (Gap Analysis):} تقوم هذه الوحدة بعمل تقييم شامل لحالة المؤسسة مقارنة بالمعيار المطلوب. تصدر تقارير تفصيلية توضح بدقة المعايير التي تم استيفاؤها، وتلك التي تفتقر إلى أدلة كافية، مع اقتراح توصيات عملية لسد هذه الفجوات. \\
4. \textbf{لوحات القيادة والتحليلات (Dashboards & Analytics):} توفر واجهات رسومية تفاعلية تعرض مؤشرات الأداء الرئيسية (KPIs) في الوقت الفعلي. تتيح للإدارة العليا وصناع القرار رؤية شاملة لمستوى التقدم، نسبة الإنجاز، والمناطق التي تحتاج إلى تدخل فوري. \\
5. \textbf{إدارة المعايير (Standards Hub):} تسمح للمؤسسات باستيراد وتخصيص أطر الاعتماد المختلفة، مما يضفي مرونة عالية للنظام ليناسب أي جهة اعتماد محلية أو دولية. \\

تم بناء النظام باستخدام تقنيات متطورة تضمن السرعة والموثوقية، حيث تم استخدام إطار عمل Next.js لتطوير واجهة المستخدم الأمامية (Front-End)، و FastAPI لتطوير الواجهة الخلفية (Back-End)، مع الاعتماد على قواعد بيانات PostgreSQL المدعومة بـ Supabase و pgvector للتعامل مع البيانات الضخمة والمتجهات الذكية. \\

\section*{1.3 تعريف المشكلة:}
\addcontentsline{toc}{section}{1.3 تعريف المشكلة:}

تعاني العديد من المؤسسات التعليمية في سعيها للحصول على الاعتماد الأكاديمي من تحديات جمة وعقبات تعرقل مسيرتها نحو التميز. يمكن تلخيص المشكلة الرئيسية في \textbf{"الافتقار إلى نظام متكامل وذكي يدير تعقيدات عملية ضمان الجودة، والاعتماد المفرط على العمليات اليدوية المشتتة"}. \\

وتتفرع من هذه المشكلة الأساسية عدة مشكلات فرعية مؤثرة: \\
\begin{itemize}
\item \textbf{العمل اليدوي المرهق والمستهلك للوقت:} تجميع الوثائق، فرزها، وتسكينها يدوياً أمام كل مؤشر أو معيار هي عملية تتطلب مئات الساعات من العمل الإداري والأكاديمي، مما يشتت جهود أعضاء هيئة التدريس والإداريين عن مهامهم الأساسية.
\item \textbf{تشتت البيانات وغياب المركزية:} غالباً ما تكون الأدلة والوثائق موزعة عبر أجهزة حواسيب مختلفة، أو مخزنة في ملفات ورقية، أو عبر منصات تخزين غير مترابطة. هذا التشتت يجعل من الصعب جداً استرجاع المعلومة المناسبة في الوقت المناسب، خاصة عند اقتراب مواعيد زيارات لجان المراجعة الخارجية.
\item \textbf{صعوبة تحديد الفجوات بشكل استباقي:} في ظل غياب أدوات التحليل الذكية، تكتشف المؤسسات أوجه القصور في ملفات الجودة الخاصة بها في أوقات متأخرة، أحياناً خلال عملية التدقيق نفسها، مما يعرضها لفقدان الاعتماد أو تأجيله. تحديد الفجوات يدوياً هو عملية عرضة للتحيز والخطأ البشري.
\item \textbf{غياب الرؤية الشاملة والمتابعة اللحظية (Lack of Real-time Insights):} لا تمتلك القيادات الأكاديمية أدوات تتيح لهم معرفة نسبة الإنجاز الفعلية في تجهيز ملفات الجودة بشكل لحظي، وتعتمد على تقارير دورية قد تكون غير محدثة أو تفتقر إلى الدقة.
\item \textbf{التحديات في التنسيق والتعاون:} عملية الاعتماد تتطلب تضافر جهود عدة أقسام (شؤون الطلاب، الموارد البشرية، الشؤون الأكاديمية، الخ). الأساليب التقليدية لا توفر بيئة تعاونية تتيح تتبع المهام ومشاركة التحديثات بشكل سلس وفعال.
\end{itemize}

من هنا جاءت فكرة منصة "Ayn" لتقدم حلاً جذرياً لهذه المشاكل، من خلال أتمتة العمليات الروتينية، توفير مركزية مطلقة للبيانات، وتسخير قدرات الذكاء الاصطناعي لتقديم تحليلات دقيقة ولحظية. \\

\section*{1.4 الأهداف:}
\addcontentsline{toc}{section}{1.4 الأهداف:}

يهدف مشروع منصة "عين" إلى إحداث ثورة في مفهوم إدارة الجودة الأكاديمية من خلال تحقيق مجموعة من الأهداف الاستراتيجية والتشغيلية التقنية: \\

\textbf{الأهداف الرئيسية (الاستراتيجية):} \\
1. \textbf{تمكين المؤسسات التعليمية:} مساعدة الجامعات والمعاهد والمدارس على تحقيق معايير الجودة الوطنية والدولية بكفاءة عالية وبأقل جهد ووقت ممكن. \\
2. \textbf{التحول الرقمي الشامل لملفات الجودة:} التخلص التدريجي من المعاملات الورقية والأنظمة المشتتة، واستبدالها ببيئة عمل سحابية آمنة وموحدة. \\
3. \textbf{تعزيز ثقافة الجودة المستدامة:} تحويل ضمان الجودة من مجرد "حدث موسمي" استعداداً لزيارات الاعتماد، إلى ممارسة يومية مستمرة وثقافة متأصلة داخل المؤسسة. \\

\textbf{الأهداف التشغيلية والتقنية:} \\
1. \textbf{أتمتة إدارة وتصنيف الأدلة:} تطوير نظام ذكي قادر على استقبال الملفات، قراءتها، وفهرستها تلقائياً باستخدام تقنيات معالجة اللغات الطبيعية. \\
2. \textbf{الكشف المبكر والدقيق عن الفجوات:} بناء خوارزميات ذكية تقوم بمطابقة الأدلة المتوفرة مع متطلبات المعايير، وتحديد أوجه النقص بدقة عالية، مما يوفر على المراجعين البشريين ساعات طويلة من المراجعة اليدوية. \\
3. \textbf{تقديم مساعد ذكي (Horus AI):} توفير واجهة تفاعلية تعتمد على الذكاء الاصطناعي التوليدي للرد على استفسارات المستخدمين حول معايير الجودة، واقتراح خطط تصحيحية بناءً على سياق المؤسسة. \\
4. \textbf{توفير تحليلات مرئية متقدمة:} تصميم لوحات تحكم (Dashboards) تفاعلية تقدم إحصائيات دقيقة، ورسوماً بيانية توضح نسب الإنجاز، وتساعد الإدارة العليا في اتخاذ قرارات مبنية على بيانات حقيقية ومحدثة لحظياً. \\
5. \textbf{ضمان أمن وسرية البيانات:} بناء هيكلية أمنية قوية تعتمد على التشفير والتحكم الدقيق في الصلاحيات (RBAC) لضمان حماية الوثائق الحساسة والبيانات المؤسسية من أي وصول غير مصرح به. \\

\section*{1.5 الدراسات السابقة:}
\addcontentsline{toc}{section}{1.5 الدراسات السابقة:}

في إطار السعي لتطوير منصة "Ayn"، تم إجراء دراسة مستفيضة للأنظمة والحلول الحالية المستخدمة في مجال إدارة الجودة الأكاديمية. وقد أظهرت المراجعة للأدبيات والأنظمة المتاحة في السوق ما يلي: \\

1. \textbf{الأنظمة التقليدية لإدارة الوثائق (DMS):} \\
تعتمد العديد من المؤسسات على أنظمة مثل Microsoft SharePoint أو Google Workspace كحلول لتخزين ومشاركة ملفات الجودة. \\
\begin{itemize}
\item \textbf{القصور:} هذه الأنظمة توفر تخزيناً مركزياً ممتازاً، ولكنها تفتقر تماماً للسياق الأكاديمي. لا تستطيع هذه الأنظمة فهم "ما هو المعيار" أو "ما إذا كان هذا الملف يلبي متطلبات الاعتماد". تظل عملية الربط والتحليل يدوية بالكامل.
\end{itemize}

2. \textbf{أنظمة تخطيط موارد المؤسسات التعليمية (ERP & SIS):} \\
تمتلك المؤسسات التعليمية أنظمة لإدارة شؤون الطلاب (SIS) والموارد البشرية. \\
\begin{itemize}
\item \textbf{القصور:} بالرغم من غناها بالبيانات، إلا أنها غير مصممة لإدارة تدفقات عمل الجودة. استخراج التقارير منها لتلبية متطلبات هيئات الاعتماد يكون غالباً معقداً ويتطلب تدخلات يدوية لتنسيقها بما يتوافق مع النماذج المطلوبة.
\end{itemize}

3. \textbf{أنظمة إدارة الجودة الأكاديمية المتخصصة الحالية:} \\
توجد بعض البرمجيات المتخصصة في السوق والتي تقدم قوالب جاهزة ونماذج لربط الأدلة بالمعايير. \\
\begin{itemize}
\item \textbf{القصور:} معظم هذه الأنظمة تعمل كقواعد بيانات علائقية تقليدية (CRUD applications). تعتمد بالكامل على المدخلات البشرية، حيث يجب على المستخدم قراءة الدليل، تقييمه، ثم ربطه يدوياً بالمعيار. كما أن الكثير منها يعاني من واجهات مستخدم معقدة وقديمة (Legacy UI).
\end{itemize}

\textbf{الفجوة التي يعالجها مشروع Ayn:} \\
تكمن الفجوة البحثية والتقنية الحقيقية في \textbf{"غياب الذكاء الاصطناعي السياقي المدمج في قلب دورة حياة الجودة"}. الأنظمة السابقة تعتبر أدوات "سلبية" تنتظر إدخال البيانات وتوجيه المستخدم. أما منصة Ayn، فهي تقدم نهجاً "استباقياً" وتفاعلياً. من خلال دمج قدرات Retrieval-Augmented Generation (RAG) ونماذج LLMs المتقدمة، لا يقوم النظام بتخزين الملف فحسب، بل يقرأ محتواه، يفهم سياقه، يقترح المعيار المناسب له، ويقيم مدى جودته. هذا التطور ينقل النظام من مجرد "أرشيف إلكتروني" إلى "مستشار جودة رقمي" يعمل على مدار الساعة. \\

\section*{1.6 خطة الدراسة:}
\addcontentsline{toc}{section}{1.6 خطة الدراسة:}

لضمان نجاح المشروع وتطوير منصة Ayn لتخرج بصورة احترافية وتلبي كافة المتطلبات التقنية والأكاديمية، تم اتباع منهجية علمية وهندسية دقيقة لأسلوب تطوير البرمجيات، وتقسيم العمل إلى عدة مراحل متسلسلة ومتداخلة (Iterative Process) تتمثل في الآتي: \\

\textbf{1. مرحلة الدراسة والتحليل (Requirement Analysis & Planning):} \\
\begin{itemize}
\item \textbf{جمع المتطلبات:} تم إجراء دراسة معمقة لمتطلبات هيئات الاعتماد الأكاديمي المحلية والدولية (مثل الهيئة القومية لضمان جودة التعليم والاعتماد بمصر، ومعايير الآيزو).
\item \textbf{تحديد الفئات المستهدفة:} تحليل أدوار المستخدمين المختلفين داخل المنصة (مدير النظام، المراجع، عضو هيئة التدريس، مسؤول الجودة) وتحديد الصلاحيات الخاصة بكل دور (Role-Based Access Control).
\item \textbf{تحليل تدفق البيانات:} رسم مخططات النظام وتحليل كيفية انتقال البيانات من لحظة رفع الملف وحتى صدور تقرير تحليل الفجوات والاعتماد النهائي.
\end{itemize}

\textbf{2. مرحلة التصميم المعماري وواجهة المستخدم (Design Phase):} \\
\begin{itemize}
\item \textbf{تصميم واجهات المستخدم (UI/UX Design):} التركيز على بناء تجربة مستخدم عصرية وبديهية (Intuitive UX) تعتمد على البساطة والفعالية، وتقليل عدد النقرات اللازمة للوصول للمعلومة، مع ضمان التوافق مع الهواتف المحمولة واللغتين العربية والإنجليزية.
\item \textbf{التصميم المعماري للنظام (System Architecture):} اختيار بنية الأنظمة السحابية الموزعة (Microservices-oriented architecture) وتحديد التكنولوجيا المناسبة (Next.js للواجهات، FastAPI للخلفية، و Prisma لإدارة قواعد البيانات).
\item \textbf{تصميم قواعد البيانات (Database Schema):} بناء هيكل بيانات مرن (ERD) يستوعب تعقيدات المعايير والمؤشرات الفرعية، ويدعم تخزين المتجهات (Vector Embeddings) اللازمة لعمليات الذكاء الاصطناعي.
\end{itemize}

\textbf{3. مرحلة التنفيذ والبرمجة (Implementation & Development):} \\
\begin{itemize}
\item \textbf{تطوير الواجهة الخلفية (Backend Development):} برمجة منطق العمل (Business Logic) المتمثل في إدارة المؤسسات، المستخدمين، المعايير، ورفع الملفات، وتطوير واجهات برمجة التطبيقات (RESTful APIs).
\item \textbf{دمج الذكاء الاصطناعي (AI Integration):} بناء وتدريب محرك Horus AI ليتفاعل مع المستندات، وتنفيذ تقنية RAG للبحث الدلالي وتوليد التحليلات التلقائية وتحديد الفجوات.
\item \textbf{تطوير الواجهة الأمامية (Frontend Development):} تحويل التصميمات إلى واجهات تفاعلية حية، وربطها مع الواجهة الخلفية لضمان عرض البيانات بشكل لحظي.
\end{itemize}

\textbf{4. مرحلة الاختبار والمراجعة (Testing & QA):} \\
\begin{itemize}
\item \textbf{اختبار الوحدات (Unit Testing):} فحص كل مكون برمجي على حدة لضمان خلوه من الأخطاء المنطقية.
\item \textbf{اختبار التكامل (Integration Testing):} التأكد من أن جميع الوحدات تتواصل مع بعضها البعض بشكل صحيح (مثل اتصال الواجهة الأمامية بالخلفية وقاعدة البيانات، واتصال محرك الذكاء الاصطناعي بالمنصة).
\item \textbf{اختبار الأداء والضغط (Performance & Load Testing):} تقييم قدرة النظام على التعامل مع رفع عدد كبير من الملفات ومعالجة البيانات المعقدة للذكاء الاصطناعي دون تباطؤ أو انهيار.
\item \textbf{اختبار قابلية الاستخدام (Usability Testing):} محاكاة سيناريوهات استخدام حقيقية لضمان سهولة تعامل المستخدمين مع المنصة وتحديد أي معوقات لتجربة المستخدم.
\end{itemize}

\textbf{5. مرحلة النشر والتطوير المستمر (Deployment & Continuous Improvement):} \\
\begin{itemize}
\item إعداد البنية التحتية السحابية (Cloud Infrastructure) ونشر التطبيق باستخدام تقنيات الحاويات (Containers) لضمان الاستقرار وسهولة التوسع المستقبلي.
\item وضع خطط مستقبلية لإضافة ميزات جديدة وتحديث نماذج الذكاء الاصطناعي لمواكبة التطورات التقنية ومتطلبات الاعتماد المتغيرة.
\end{itemize}

\newpage
\chapter*{الفصل الثاني: فكرة المشروع (Project Idea)}
\addcontentsline{toc}{chapter}{الفصل الثاني: فكرة المشروع (Project Idea)}

تعتبر فكرة مشروع منصة "عين" (Ayn) نتاجاً لدراسة عميقة لاحتياجات سوق العمل التعليمي، والوقوف على أبرز التحديات التي تواجه مسؤولي الجودة والاعتماد في المؤسسات الأكاديمية. إن الفكرة لم تُولد من فراغ، بل تبلورت كحل تقني طموح يهدف إلى القضاء على البيروقراطية الإدارية والورقية المعقدة، واستبدالها بنظام سحابي ذكي ومؤتمت. في هذا الفصل، سنسلط الضوء بشكل موسع على الأسس التي بنيت عليها فكرة المشروع، بدءاً من أهدافه الاستراتيجية والتشغيلية، مروراً بنطاقه ومحدداته، وصولاً إلى أهميته البالغة والجمهور المستهدف من وراء هذا الابتكار التكنولوجي. \\

\begin{figure}[H]
\centering
\includegraphics[width=0.8\textwidth]{extracted_images/img_02.png}
\caption{المفهوم الأساسي للنظام}
\end{figure}

\section*{2.1 أهداف المشروع (Project Objectives)}
\addcontentsline{toc}{section}{2.1 أهداف المشروع (Project Objectives)}

لقد صُممت منصة "Ayn" لتحقيق مجموعة واسعة ومترابطة من الأهداف التي تسعى في مجملها إلى تمكين المؤسسات التعليمية من بناء نظام داخلي قوي لضمان الجودة. تنقسم هذه الأهداف إلى أهداف تقنية، وأخرى إدارية وتشغيلية: \\

\subsection*{2.1.1 الأهداف التقنية والابتكارية}
\addcontentsline{toc}{subsection}{2.1.1 الأهداف التقنية والابتكارية}
1. \textbf{تطوير بنية تحتية سحابية مرنة ومتدرجة (Scalable Cloud Architecture):} بناء نظام قادر على استيعاب ومعالجة كميات هائلة من البيانات، بدءاً من ملفات نصية بسيطة وصولاً إلى ملفات PDF ضخمة تضم مئات الصفحات، مع ضمان سرعة الوصول وسلاسة الاستخدام (High Availability) حتى في أوقات ذروة الاستخدام. \\
2. \textbf{الاستفادة القصوى من قدرات الذكاء الاصطناعي التوليدي (Generative AI):} دمج نماذج لغوية متقدمة، وبالتحديد نموذج Gemini 1.5 Pro، لإنشاء المساعد الذكي "Horus AI". الهدف هو تحويل النظام من مجرد مستودع سلبي للبيانات إلى أداة تحليلية ذكية قادرة على الفهم الدلالي للوثائق (Semantic Understanding) وتقديم التوصيات. \\
3. \textbf{أتمتة دورة حياة تحليل الفجوات (Automated Gap Analysis Lifecycle):} بدلاً من الفحص اليدوي المجهد للمستندات، يهدف النظام إلى بناء خوارزمية قوية تأخذ المعيار الأكاديمي كمدخل، وتطابقه تلقائياً مع الأدلة المتاحة، ثم تصدر تقريراً تفصيلياً يوضح مواطن القوة ونقاط الضعف بكل دقة وسرعة. \\
4. \textbf{تطبيق تقنيات البحث الدلالي المتقدم (RAG - Retrieval-Augmented Generation):} إنشاء قاعدة بيانات متجهية (Vector Database) باستخدام تقنية `pgvector`، وذلك للسماح بالبحث في محتوى الملفات والأدلة ليس فقط بالكلمات المفتاحية البسيطة، ولكن بناءً على المعنى والسياق، مما يعزز من قدرة محرك الذكاء الاصطناعي على استخراج إجابات دقيقة وموثقة. \\

\subsection*{2.1.2 الأهداف الإدارية والتشغيلية}
\addcontentsline{toc}{subsection}{2.1.2 الأهداف الإدارية والتشغيلية}
1. \textbf{القضاء على المركزية وتسهيل التعاون المشترك (Enhancing Collaboration):} توفير بيئة عمل افتراضية تُمكّن مختلف الأقسام والكليات والموظفين داخل المؤسسة من العمل معاً بشكل متزامن. حيث يمكن لعضو هيئة التدريس رفع الأدلة، ليقوم مسؤول الجودة بمراجعتها، بينما تتابع الإدارة العليا تقدم العمل عبر لوحات قيادة مخصصة. \\
2. \textbf{تعزيز الشفافية وحوكمة البيانات:} ضمان تسجيل وتوثيق كل إجراء يتم داخل المنصة عبر وحدة "سجل النشاط" (Activity Logs)، مما يوفر مسار تدقيق (Audit Trail) كامل يوضح من قام برفع الدليل، ومتى، وكيف تم تقييمه، لضمان أعلى معايير الشفافية. \\
3. \textbf{دعم اتخاذ القرار المبني على البيانات (Data-Driven Decision Making):} تمكين قيادات المؤسسات التعليمية من اتخاذ قرارات استراتيجية بناءً على بيانات محدثة لحظياً تُعرض من خلال لوحات القيادة (Dashboards) التفاعلية، بدلاً من الاعتماد على تقارير ورقية دورية قد تكون قديمة أو غير دقيقة. \\
4. \textbf{توفير الوقت وخفض التكاليف التشغيلية:} من خلال أتمتة العمليات الروتينية، يتم توفير آلاف الساعات من العمل الإداري الذي كان يُهدر في تجميع الورق وطباعته وتصنيفه، مما يترجم بشكل مباشر إلى توفير هائل في التكاليف التشغيلية للمؤسسة. \\

\begin{figure}[H]
\centering
\includegraphics[width=0.8\textwidth]{extracted_images/img_03.png}
\caption{نظرة على أهداف وأقسام المنصة}
\end{figure}

\section*{2.2 نطاق وقيود المشروع (Project Scope and Limitations)}
\addcontentsline{toc}{section}{2.2 نطاق وقيود المشروع (Project Scope and Limitations)}

إن تحديد نطاق العمل بوضوح هو خطوة حاسمة لنجاح أي مشروع برمجي ضخم كمنصة Ayn. يوضح هذا القسم حدود إمكانيات المنصة في مرحلتها الحالية، بالإضافة إلى القيود الفنية أو العملية التي تم أخذها بعين الاعتبار. \\

\subsection*{2.2.1 نطاق المشروع (In Scope)}
\addcontentsline{toc}{subsection}{2.2.1 نطاق المشروع (In Scope)}
\begin{itemize}
\item \textbf{إدارة دورة حياة الأدلة:} يشمل النظام عملية رفع الملفات (PDFs, DOCX, Images)، استخراج النصوص منها، فهرستها ذكياً، وربطها بمعايير الاعتماد المختلفة.
\item \textbf{تخصيص المعايير:} المنصة مصممة لتكون مرنة، وتتيح بناء هياكل هرمية لأي نوع من أنواع الاعتماد (سواء كان يحتوي على معايير، مؤشرات، وأمثلة تطبيقية) مثل معايير ISO و NCAAA.
\item \textbf{التحليل الذكي وإنشاء التقارير:} النظام يتولى مهمة مطابقة الأدلة المرفوعة وإصدار تقارير تحليل الفجوات (Gap Analysis Reports) مع تقديم توصيات تصحيحية.
\item \textbf{لوحات التحكم والإشعارات اللحظية:} توفير نظام متكامل للإحصائيات ونظام تنبيهات فوري لضمان سرعة التجاوب بين فريق العمل.
\item \textbf{المساعد الافتراضي Horus AI:} توفير نافذة محادثة ذكية متخصصة حصراً في نطاق الجودة المؤسسية والاعتماد، تستطيع استرجاع بيانات المنصة والرد على استفسارات المراجعين.
\end{itemize}

\subsection*{2.2.2 قيود ومحددات المشروع (Limitations & Out of Scope)}
\addcontentsline{toc}{subsection}{2.2.2 قيود ومحددات المشروع (Limitations & Out of Scope)}
\begin{itemize}
\item \textbf{عدم اتخاذ قرارات مصيرية تلقائية:} الذكاء الاصطناعي في منصة Ayn (Horus AI) مصمم ليعمل بصفة "استشارية" بحتة. هو يحلل ويقترح، لكنه \textbf{لا يتخذ أي إجراءات لتغيير البيانات أو مسحها أو اعتمادها بشكل نهائي} نيابة عن المستخدم. القرار النهائي دائماً بيد العنصر البشري (مسؤول الجودة).
\item \textbf{الاعتماد على جودة المدخلات:} تعتمد دقة محرك الذكاء الاصطناعي بشكل كبير على جودة الملفات والأدلة المرفوعة (Garbage in, Garbage out). إذا تم رفع وثائق ممسوحة ضوئياً بجودة رديئة جداً غير قابلة للتعرف على الحروف (OCR)، فإن قدرة النظام على التحليل الدقيق ستتأثر.
\item \textbf{الاستهلاك العالي للموارد السحابية (Cost Restraints):} عمليات التحليل الدلالي وإنشاء الـ Vector Embeddings تتطلب موارد حوسبية مكلفة (API Tokens). تم تصميم النظام بتقنيات التخزين المؤقت (Caching) ومحدودية الاستعلامات (Rate Limiting) لتجنب الإفراط في التكاليف السحابية.
\item \textbf{التكامل المعقد مع الأنظمة القديمة (Legacy Systems):} في نسخته الحالية، لا يتضمن المشروع ربطاً آلياً مباشراً مع أنظمة تخطيط الموارد التعليمية (ERP) أو منصات التعلم (LMS) الخاصة بالجامعات لجلب البيانات منها تلقائياً، وتعتمد المنصة على قيام المستخدمين برفع المخرجات المطلوبة.
\end{itemize}

\section*{2.3 أهمية المشروع (Why it Matters)}
\addcontentsline{toc}{section}{2.3 أهمية المشروع (Why it Matters)}

في المشهد التعليمي الحديث، الذي يتسم بتنافسية شديدة، لم يعد الاعتماد الأكاديمي والحصول على شهادات الجودة مجرد إجراء روتيني اختياري، بل أصبح ضرورة حتمية ومؤشراً رئيسياً يعكس كفاءة المؤسسة وسمعتها، ويحدد مدى أهليتها للحصول على التمويل وجذب الطلاب. من هنا، تكتسب منصة "عين" أهميتها القصوى من خلال تقديم قيمة مضافة لا غنى عنها للمؤسسات. \\

أولاً، تقوم المنصة \textbf{بتخفيف العبء الإداري الهائل} المصاحب لعمليات الجودة. ففي العادة، تتطلب الاستعدادات لزيارات لجان الاعتماد أشهراً من التجهيز، وجمع الملفات، وعقد الاجتماعات الطويلة لمطابقة الوثائق مع المؤشرات. "عين" تحول هذا الكابوس الإداري إلى عملية سلسة، رقمية، ومنظمة تتيح للموظفين التركيز على الجوهر (الارتقاء بالعملية التعليمية) بدلاً من الانشغال بالإجراءات المكتبية. \\

ثانياً، تمنح المنصة المؤسسات \textbf{الثقة واليقين} قبل الزيارات الرسمية لهيئات الاعتماد. من خلال الرؤى اللحظية والتحليل العميق الذي يقدمه المساعد الذكي Horus AI، تستطيع الإدارة تحديد نقاط ضعفها ومعالجتها مبكراً (Proactive Remediation)، مما يقلل بشكل كبير من مخاطر الفشل في الحصول على الاعتماد أو تلقي ملاحظات سلبية من المدققين الخارجيين. \\

ثالثاً، تساهم المنصة في \textbf{تحويل ثقافة المؤسسة}. فبدلاً من أن تكون الجودة مسؤولية شخص واحد أو قسم بعينه، تتيح أدوات التعاون في المنصة إشراك جميع الأطراف المعنية (هيئة تدريس، إداريين، قيادات) في عملية ضمان الجودة، لتصبح ممارسة مؤسسية شاملة ومستدامة. \\

\section*{2.4 المستخدمون المستهدفون والمؤسسات (Target Users and Institutions)}
\addcontentsline{toc}{section}{2.4 المستخدمون المستهدفون والمؤسسات (Target Users and Institutions)}

تم بناء هيكلية المنصة بحيث تكون مرنة وقابلة للتخصيص (Multi-tenant Architecture) لتخدم شريحة واسعة ومتنوعة من الكيانات الأكاديمية والمهنية، والمستخدمين ذوي الصلاحيات المتدرجة. \\

\subsection*{2.4.1 المؤسسات المستهدفة}
\addcontentsline{toc}{subsection}{2.4.1 المؤسسات المستهدفة}
\begin{itemize}
\item \textbf{الجامعات والكليات:} المؤسسات الكبرى ذات الهياكل الإدارية المعقدة، والتي تحتاج لإدارة ملفات الجودة عبر عدة كليات وأقسام أكاديمية مختلفة تحت مظلة اعتماد واحدة أو عدة اعتمادات.
\item \textbf{المعاهد العليا ومراكز التدريب:} التي تسعى لتطبيق معايير الاعتماد المهني أو الأكاديمي لضمان الاعتراف بمخرجاتها التعليمية في سوق العمل.
\item \textbf{المدارس الحكومية والخاصة:} التي تهدف للامتثال لمعايير ضمان الجودة المحلية (مثل معايير الجودة والاعتماد المصرية NAQAAE).
\item \textbf{هيئات ومنظمات الاعتماد نفسها:} يمكن تكييف المنصة مستقبلاً لتكون أداة تستخدمها اللجان المانحة للاعتماد في استقبال ملفات المتقدمين، مراجعتها، وتنظيم زيارات المراجعة.
\end{itemize}

\subsection*{2.4.2 أدوار المستخدمين المستهدفين داخل النظام}
\addcontentsline{toc}{subsection}{2.4.2 أدوار المستخدمين المستهدفين داخل النظام}
ولأن النظام يخدم مؤسسة بأكملها، تم تصميم عدة أدوار بصلاحيات مختلفة (RBAC) لضمان أمان البيانات وتقسيم المهام: \\
\begin{itemize}
\item \textbf{مدير النظام (Admin / Institution Manager):} يتمتع بصلاحيات شاملة لتهيئة النظام، إدارة حسابات المستخدمين، واستيراد أو تعديل أطر الاعتماد والمعايير الرئيسية.
\item \textbf{مسؤول الجودة / المراجع الداخلي (Auditor):} المحرك الرئيسي للنظام. يقوم بمراجعة الأدلة المرفوعة، إجراء تحليلات الفجوات، والتواصل مع محرك الذكاء الاصطناعي Horus AI لفهم مدى مطابقة الأدلة، واعتماد خطط التحسين.
\item \textbf{عضو هيئة التدريس / الموظف (Teacher / Staff):} دور يقتصر على تنفيذ المهام الموكلة إليه، مثل توفير ورفع الأدلة المحددة، والرد على ملاحظات مسؤولي الجودة دون صلاحية التعديل على هياكل المعايير الأساسية.
\item \textbf{الإدارة العليا / صانع القرار (Management / University Role):} حسابات ذات طبيعة إشرافية، تركز بالأساس على لوحات القيادة (Dashboards) والتقارير الإحصائية (Analytics) لمتابعة الأداء العام دون الغوص في التفاصيل التشغيلية.
\end{itemize}

\section*{2.5 القيمة الأساسية المقترحة (Core Value Proposition)}
\addcontentsline{toc}{section}{2.5 القيمة الأساسية المقترحة (Core Value Proposition)}

تتمثل القيمة الجوهرية (Core Value Proposition) التي تقدمها منصة Ayn في تحويل مفهوم ضمان الجودة الأكاديمية من \textbf{"عبء إداري وتجميع للمستندات"} إلى \textbf{"مسعى ذكي، كفؤ، واستراتيجي للتطوير المستمر"}. هذا التحول الجذري يتم تحقيقه من خلال حزمة من المزايا الحصرية التي تميز النظام وتجعله أداة استثنائية في السوق: \\

\begin{longtable}{|c|c|}
\hline
الميزة التنافسية (Feature) & الوصف الممتد للقيمة المضافة (Detailed Description) \\ \hline
:--- & :--- \\ \hline
\textbf{الامتثال المدعوم بالذكاء الاصطناعي (AI-Powered Compliance)} & لا تعتمد المنصة على الجهد البشري وحده في قراءة آلاف الصفحات. محرك Horus AI يقوم بأتمتة تصنيف الأدلة، استكشاف الفجوات بشكل استباقي، واقتراح خطط تصحيحية ذكية، مما يضاعف سرعة العمل ودقته بشكل غير مسبوق. \\ \hline
\textbf{بيئة عمل سحابية وموحدة (Unified Single Source of Truth)} & توفير مصدر واحد وموثوق (Single Source of Truth) يجمع كل ما يخص الجودة؛ من الوثائق، إلى نصوص المعايير، وسير العمل. هذا يقضي تماماً على تشتت النسخ والملفات الضائعة بين الأقسام. \\ \hline
\textbf{رؤى لحظية دقيقة (Real-time Actionable Insights)} & عبر لوحات التحكم والمخططات التفاعلية، تتحول البيانات الجامدة إلى رؤى حية. يستطيع صناع القرار بلمحة واحدة معرفة الموقف الحالي للامتثال وتحديد نقاط الخطر مبكراً لاتخاذ إجراءات فورية. \\ \hline
\textbf{تبسيط وتنسيق سير العمل (Streamlined Workflows)} & النظام يوفر بيئة توجيهية (Guided Workflows) تنقل المستخدم خطوة بخطوة؛ من استيراد المعيار، مروراً برفع الأدلة، ثم طلب التحليل، وأخيراً استعراض التقرير. مما يخفف من التخبط العشوائي في المهام. \\ \hline
\textbf{الاستدامة والاعتمادية (Scalability & Adaptability)} & المنصة صممت لتبدأ مع مؤسسة صغيرة ببرنامج واحد، ولها القدرة التقنية على التوسع لتشمل جامعة بأكملها بعشرات الكليات وآلاف المستخدمين، مع إمكانية استيراد أي معيار جودة جديد بكل سهولة ويسر. \\ \hline
\end{longtable}

\begin{figure}[H]
\centering
\includegraphics[width=0.8\textwidth]{extracted_images/img_05.png}
\caption{Horus AI Assistant}
\end{figure}

خلاصة القول، إن منصة "عين" لا تقدم مجرد أداة برمجية عادية، بل هي تقدم \textbf{"الذكاء والأدوات الضرورية"} التي تسلح المؤسسات التعليمية لمواجهة تعقيدات الاعتماد الأكاديمي بكل ثقة، محولة إياه من تحدٍ يبعث على التوتر والقلق، إلى ميزة تنافسية تُبرز ريادة المؤسسة وتميزها التعليمي. \\

\newpage
\chapter*{الفصل الثالث: الأدوات والتقنيات (Tools & Technologies)}
\addcontentsline{toc}{chapter}{الفصل الثالث: الأدوات والتقنيات (Tools & Technologies)}

يعتبر اختيار التكنولوجيا المناسبة لتطوير منصة بحجم "Ayn" قراراً هندسياً حاسماً، حيث يتطلب النظام قدرة عالية على التعامل مع البيانات الضخمة، توفير أداء سريع واستجابة لحظية، وضمان مرونة عالية للدمج مع أدوات الذكاء الاصطناعي. في هذا الفصل، نستعرض بشكل مفصل ومنهجي حزمة التقنيات (Tech Stack) التي تم تبنيها لبناء هذا النظام، وتبرير أسباب اختيار كل تقنية بناءً على أهداف المشروع، بالإضافة إلى توضيح كيفية تفاعلها مع بعضها البعض لتكوين بنية تحتية متينة وقابلة للتوسع. \\

\section*{3.1 التقنيات المستخدمة في الواجهة الأمامية (Frontend Technologies)}
\addcontentsline{toc}{section}{3.1 التقنيات المستخدمة في الواجهة الأمامية (Frontend Technologies)}

تم تصميم واجهة المستخدم في منصة Ayn لتكون تطبيق صفحة واحدة (Single Page Application - SPA) يوفر تجربة تفاعلية تشبه تطبيقات سطح المكتب الأصيلة، مع الحفاظ على سرعة التحميل والوصول عبر الويب. لضمان هذه التجربة العالية الجودة، تم الاعتماد على أحدث تقنيات الويب: \\

1. \textbf{Next.js (App Router):} إطار العمل الأساسي للواجهة الأمامية. يوفر Next.js قدرات مذهلة في معالجة الصفحات سواء عن طريق التقديم من جانب الخادم (Server-Side Rendering) أو التوليد الثابت (Static Site Generation)، مما يضمن سرعة استجابة فائقة للمستخدمين وتجربة تنقل سلسة للغاية دون الحاجة لتحميل الصفحة من جديد. كما يسهم إطار العمل في توفير بنية تحتية قوية لإدارة المسارات وتسهيل عمليات تحسين محركات البحث (SEO) في المستقبل إن لزم الأمر. \\
2. \textbf{React 19:} مكتبة بناء واجهات المستخدم الشهيرة التي تعتمد على المكونات (Components). يوفر React إمكانية بناء أجزاء من الواجهة يمكن إعادة استخدامها في جميع أنحاء التطبيق، مما يُسهل من صيانة الكود البرمجي ويقلل من تكراره. \\
3. \textbf{TypeScript:} لغة برمجة قوية من Microsoft تضيف النوع الثابت (Static Typing) إلى JavaScript. استخدام TypeScript قلل بشكل كبير من الأخطاء البرمجية أثناء التطوير وضمان مستوى أعلى من الموثوقية واستقرار الكود، خاصة في مشروع معقد يحتوي على هياكل بيانات كبيرة. \\
4. \textbf{Tailwind CSS:} إطار عمل لتصميم صفحات الويب يعتمد على الأدوات (Utility-first). أتاح Tailwind CSS تصميم واجهات عصرية وجذابة بشكل سريع جداً دون الحاجة لكتابة أكواد CSS معقدة ومخصصة، مما حافظ على تجانس التصميم في كافة أرجاء المنصة وسرعة الأداء في المتصفحات. \\
5. \textbf{Shadcn UI & Radix UI:} مكتبات مكونات واجهة المستخدم التي توفر عناصر جاهزة، مصممة بشكل احترافي ومتوافقة مع معايير الوصول (Accessibility). ساهمت هذه المكتبات في تسريع عملية تصميم العناصر التفاعلية كالقوائم المنسدلة، والنوافذ المنبثقة، والأزرار. \\
6. \textbf{Zustand & React Context API:} لإدارة حالة التطبيق (State Management). تم استخدام Zustand لإدارة الحالات المعقدة بطريقة خفيفة وفعالة، بينما استُخدم الـ Context API للحالات العامة مثل حالة المصادقة وثيم التطبيق (Dark/Light Mode). \\
7. \textbf{Framer Motion:} مكتبة رسوم متحركة قوية تم استخدامها لإضافة لمسات حركية وانتقالات سلسة بين الشاشات والعناصر، مما يزيد من جاذبية التطبيق وحيوية تجربة المستخدم. \\

\begin{figure}[H]
\centering
\includegraphics[width=0.8\textwidth]{extracted_images/img_06.png}
\caption{تقنيات الواجهة الأمامية}
\end{figure}

\section*{3.2 التقنيات المستخدمة في الواجهة الخلفية (Backend Technologies)}
\addcontentsline{toc}{section}{3.2 التقنيات المستخدمة في الواجهة الخلفية (Backend Technologies)}

بينما تهتم الواجهة الأمامية بتجربة المستخدم، فإن الواجهة الخلفية هي "العمود الفقري" الذي يدير العمليات المنطقية، قواعد البيانات، والعمليات الثقيلة. صُممت الواجهة الخلفية لمنصة Ayn لتكون فائقة الأداء، قابلة للتطوير (Scalable)، وتعتمد بشكل رئيسي على بيئة بايثون (Python). \\

1. \textbf{FastAPI (Python 3.12):} إطار عمل عصري عالي الأداء لبناء واجهات برمجة التطبيقات (APIs). يتميز FastAPI بسرعته الفائقة، واعتماده الكامل على الأكواد غير المتزامنة (Asynchronous I/O) من خلال مكتبة `asyncio`، مما يجعله قادراً على معالجة آلاف الطلبات في وقت واحد دون توقف. كما يوفر توثيقاً تلقائياً وتفاعلياً للـ API باستخدام Swagger UI. \\
2. \textbf{Uvicorn:} خادم ويب عالي الأداء من نوع ASGI مخصص لتطبيقات الويب غير المتزامنة في بايثون. يعمل Uvicorn كمحرك تشغيل لتطبيق FastAPI لضمان أعلى درجات الأداء. \\
3. \textbf{Prisma Client Python:} أداة متطورة للتعامل مع قواعد البيانات (ORM). بدلاً من كتابة استعلامات SQL معقدة وعرضة للخطأ، يوفر Prisma طبقة آمنة ومكتوبة النوع (Type-safe) للتفاعل مع قاعدة البيانات، بالإضافة إلى إدارة هيكل قاعدة البيانات وتحديثاتها (Migrations) بكل سهولة. \\
4. \textbf{Pydantic:} مكتبة أساسية في بايثون لضمان صحة ونوع البيانات المدخلة والمخرجة. يتم استخدام Pydantic لإنشاء نماذج البيانات (Models) والتأكد من مطابقتها للشروط المحددة قبل معالجتها أو حفظها في قاعدة البيانات. \\
5. \textbf{مصادقة JWT (JSON Web Tokens):} نظام قوي وآمن لمصادقة المستخدمين بدون الحاجة لتخزين جلسات على الخادم (Stateless Authentication). بمجرد تسجيل الدخول بنجاح، يحصل المستخدم على رمز مميز مشفر (Token) يستخدم للتحقق من هويته في كافة الطلبات اللاحقة. \\

\begin{figure}[H]
\centering
\includegraphics[width=0.8\textwidth]{extracted_images/img_07.png}
\caption{تقنيات الواجهة الخلفية}
\end{figure}

\section*{3.3 قواعد البيانات والتخزين المؤقت (Databases & Caching)}
\addcontentsline{toc}{section}{3.3 قواعد البيانات والتخزين المؤقت (Databases & Caching)}

معمارية البيانات في منصة Ayn تتطلب نظاماً قادراً على التعامل مع بيانات علائقية معقدة (مثل العلاقات بين المؤسسات، المعايير، والمؤشرات)، بالإضافة إلى قدرة استثنائية على تخزين المتجهات الذكية (Vectors) لتمكين البحث الدلالي. \\

1. \textbf{PostgreSQL:} نظام إدارة قواعد البيانات العلائقية المتقدم والأكثر موثوقية في العالم مفتوح المصدر. تم اختياره لقدرته الهائلة على التوسع، أمانه، ودعمه الممتاز للاستعلامات المعقدة. \\
2. \textbf{Supabase:} منصة شاملة تم توظيفها كبديل مفتوح المصدر لـ Firebase. تقدم Supabase قاعدة بيانات PostgreSQL مُدارة بالكامل، بالإضافة إلى خدمة تخزين الملفات السحابية (Supabase Storage) والتي تم استخدامها لتخزين كافة المستندات المرفوعة من قبل المستخدمين (PDFs, Images) بشكل آمن ومنظم. \\
3. \textbf{مكتية pgvector:} إضافة (Extension) جوهرية لقاعدة بيانات PostgreSQL تتيح تخزين والبحث في الـ Vector Embeddings. بدون هذه التقنية، لم يكن من الممكن لمحرك الذكاء الاصطناعي "Horus AI" إجراء عمليات البحث الدلالي المعقدة واسترجاع المعلومات بناءً على السياق وليس الكلمات المفتاحية فقط (RAG). \\
4. \textbf{Upstash Redis:} نظام تخزين متطور في الذاكرة المؤقتة (In-memory Caching). تم استخدامه بشكل أساسي لتخزين نتائج الاستعلامات المتكررة والإحصائيات الخاصة بلوحات القيادة (Dashboards). هذا الإجراء يقلل الضغط بشكل هائل على قاعدة البيانات الرئيسية ويضمن تحميل البيانات على واجهة المستخدم في أجزاء من الثانية. \\

\section*{3.4 دمج الذكاء الاصطناعي (AI Integration)}
\addcontentsline{toc}{section}{3.4 دمج الذكاء الاصطناعي (AI Integration)}

لتحقيق الرؤية الأساسية للمنصة كمساعد ذكي، كان لا بد من دمج تقنيات ونماذج لغوية قوية ضمن البنية التحتية للنظام: \\

1. \textbf{Google Gemini 1.5 Pro:} هو المحرك الرئيسي للذكاء الاصطناعي في المنصة. تم اختيار هذا النموذج بالذات لقدرته الفائقة على معالجة سياقات طويلة جداً (Long Context Window)، مما يجعله مثالياً لقراءة آلاف الصفحات من المستندات المؤسسية وتحليلها دفعة واحدة. \\
2. \textbf{OpenRouter (Fallback mechanism):} لضمان استمرارية الخدمة (High Availability) في حال توقف أو بطء استجابة مزود الذكاء الاصطناعي الرئيسي (Gemini)، تم دمج OpenRouter كمزود بديل ومرن (Fallback). من خلال OpenRouter، يمكن للنظام التبديل تلقائياً لاستخدام نماذج ذكاء اصطناعي أخرى من مزودين مختلفين للحفاظ على عمل المنصة دون أي انقطاع يشعر به المستخدم. \\
3. \textbf{OpenAI API (Whisper):} تم استخدام تقنيات OpenAI بشكل متخصص في عمليات تحويل الصوت إلى نص (Speech-to-Text)، لتقديم خيارات إدخال مرنة للمستخدمين وربما لدعم الميزات المستقبلية للمحادثات الصوتية. \\

\begin{figure}[H]
\centering
\includegraphics[width=0.8\textwidth]{extracted_images/img_08.png}
\caption{خريطة الذكاء الاصطناعي والتكامل}
\end{figure}

\section*{3.5 النشر والتكامل المستمر (Deployment & DevOps)}
\addcontentsline{toc}{section}{3.5 النشر والتكامل المستمر (Deployment & DevOps)}

العمليات التشغيلية للمنصة تعتمد على بنية تحتية سحابية متطورة لتسهيل التحديثات وضمان استقرار الخدمة: \\

\begin{itemize}
\item \textbf{Vercel:} تم نشر واجهة المستخدم الأمامية (Next.js) بالكامل على منصة Vercel التي تقدم استضافة فائقة السرعة مع قدرات جبارة في التوسع التلقائي وتوزيع المحتوى عبر شبكات CDN العالمية، مما يضمن وصولاً سريعاً من أي مكان في العالم.
\item \textbf{Docker:} لم يتم إغفال أهمية تشغيل التطبيق في بيئات معزولة وموحدة. يمكن استخدام Docker لحزم الواجهة الخلفية بكل متطلباتها البرمجية لتسهيل نقلها وتشغيلها على أي خادم سحابي بسلاسة تامة.
\item \textbf{Git & GitHub:} أدوات التحكم في الإصدارات والتعاون الأساسية، حيث تُستخدم لإدارة الأكواد المصدرية، تتبع التغييرات، وتسهيل العمل الجماعي بين أعضاء فريق التطوير بكل أمان وفعالية.
\end{itemize}

خلاصة القول، إن البنية التقنية لمنصة Ayn لم تُبنَ بمحض الصدفة، بل تم اختيار كل عنصر فيها بعناية فائقة لضمان تناغمه مع باقي أجزاء النظام، وليكوّنوا معاً بيئة تقنية قوية، حديثة، مستقرة، قادرة على استيعاب تطلعات المؤسسات التعليمية وتوسعاتها المستقبلية دون أي عقبات تقنية. \\

\newpage
\chapter*{الفصل الرابع: تحليل النظام والهندسة المعمارية (System Analysis & Architecture)}
\addcontentsline{toc}{chapter}{الفصل الرابع: تحليل النظام والهندسة المعمارية (System Analysis & Architecture)}

النجاح في بناء أنظمة معقدة ومتقدمة تكنولوجياً مثل منصة "عين" لا يعتمد فقط على كتابة الأكواد، بل يتطلب في الأساس تحليلاً دقيقاً للمتطلبات وتصميماً هندسياً معمارياً متيناً يضمن سلاسة تدفق البيانات وقابلية النظام للتوسع والصيانة. في هذا الفصل، نستعرض بالتفصيل المعمارية البرمجية للمنصة، وهيكلة قواعد البيانات، ونموذج تدفق البيانات الذي يضمن الأمان والسرعة. \\

\begin{figure}[H]
\centering
\includegraphics[width=0.8\textwidth]{extracted_images/img_09.png}
\caption{المعمارية العامة للنظام}
\end{figure}

\section*{4.1 وحدات المنصة وهيكليتها (Platform Modules)}
\addcontentsline{toc}{section}{4.1 وحدات المنصة وهيكليتها (Platform Modules)}

اعتمدنا في تصميم النظام على مفهوم التصميم المبني على المجالات أو الوحدات المنفصلة (Domain-Driven Design). هذا يعني أن النظام لا يتكون من كتلة برمجية واحدة متشابكة، بل يتألف من عدة وحدات برمجية (Modules) مستقلة، كل وحدة تؤدي غرضاً محدداً وتتواصل مع الوحدات الأخرى عبر واجهات برمجية واضحة. هذا النهج يضمن أنه في حالة حدوث عطل في وحدة معينة، لا تتأثر باقي وحدات النظام بالكامل. \\

\textbf{تشمل الوحدات الأساسية للمنصة ما يلي:} \\
1. \textbf{وحدة المصادقة وإدارة المستخدمين (Authentication & Users Module):} تتحكم في عمليات تسجيل الدخول، إنشاء الحسابات، توليد وإدارة الرموز المشفرة (JWT)، والتحكم في صلاحيات الوصول بناءً على أدوار المستخدمين (RBAC). \\
2. \textbf{وحدة إدارة المؤسسات (Institutions Module):} مختصة بإدارة بيانات الكيانات الأكاديمية والمؤسسية المشتركة في المنصة، بالإضافة إلى إدارة الاشتراكات وربط الموظفين والمستخدمين بالمؤسسات التابعين لها (Multi-tenancy). \\
3. \textbf{وحدة إدارة المعايير (Standards Module):} تمثل المكتبة المركزية لجميع أطر الجودة والاعتماد. تدير هذه الوحدة استيراد المعايير، وتجزئتها إلى معايير فرعية ومؤشرات قابلة للقياس، وربط هذه المعايير بالمؤسسات التي تود التقييم بناءً عليها. \\
4. \textbf{وحدة مستودع الأدلة (Evidence Vault Module):} مسؤولة عن عمليات الرفع (Upload) والتخزين الآمن للملفات على Supabase، ومتابعة حالات معالجة الملفات (Processing States) قبل وبعد تحليلها بواسطة الذكاء الاصطناعي. \\
5. \textbf{وحدة تحليل الفجوات (Gap Analysis Module):} تعتبر من أعقد وحدات النظام؛ حيث تقوم بتجميع البيانات من وحدة الأدلة ووحدة المعايير، وإجراء المطابقات المنطقية، لتحديد المؤشرات غير المستوفاة وإصدار التقارير النهائية بنسب الإنجاز. \\
6. \textbf{وحدة الإشعارات والنشاط (Notifications & Activity Module):} نظام المراسلة الداخلي الذي يرسل تنبيهات لحظية للمستخدمين عند اكتمال عمليات المعالجة الطويلة للذكاء الاصطناعي، كما يحتفظ بسجل تاريخي (Audit Trail) لكل الحركات على المنصة. \\
7. \textbf{محرك Horus AI المركزي:} وهو المحرك الذي يتصل بـ Google Gemini ويستخدم تقنيات الـ RAG للإجابة على التساؤلات وإجراء التحليلات الدلالية بذكاء بناءً على السياق المؤسسي. \\

\section*{4.2 تدفق البيانات (Data Flow Architecture)}
\addcontentsline{toc}{section}{4.2 تدفق البيانات (Data Flow Architecture)}

يُعد تدفق البيانات من واجهة المستخدم وصولاً إلى قاعدة البيانات ثم إلى مزود الذكاء الاصطناعي والعكس، عملية حساسة تتطلب سرعة وأماناً. تم تصميم تدفق العمل (Workflow) لضمان عدم توقف واجهة المستخدم (Non-blocking UI) أثناء العمليات الحسابية المعقدة. \\

\textbf{مسار تدفق البيانات عند رفع وثيقة جودة:} \\
1. \textbf{استقبال الطلب (Request Initiation):} يختار المستخدم الملف (PDF مثلاً) ويضغط على زر الرفع من الواجهة الأمامية. \\
2. \textbf{نقل وتخزين الملف (Storage Layer):} تقوم الواجهة الخلفية باستقبال الملف ورفعه مباشرة وسريعاً إلى وحدة التخزين السحابية (Supabase Storage). بعد نجاح الرفع، يتم إنشاء سجل أولي للملف في قاعدة بيانات PostgreSQL بحالة "قيد المعالجة" (Processing). \\
3. \textbf{المعالجة الخلفية غير المتزامنة (Asynchronous Background Task):} هنا يأتي دور الـ FastAPI؛ بدلاً من جعل المستخدم ينتظر، يُرسل الخادم رداً فورياً بنجاح عملية الرفع (Http 200 OK)، ويقوم بإحالة عملية تحليل الملف لمحرك في الخلفية (Background Task). \\
4. \textbf{التحليل الذكي (AI Analysis & Embedding):} يقوم محرك Horus AI في الخلفية باستخراج النصوص من الملف المرفوع، وتقسيمها إلى نصوص أصغر (Chunking)، ثم تحويلها إلى متجهات رياضية (Vectors) عبر واجهات الـ AI. \\
5. \textbf{تخزين المتجهات (Vector Storage):} تُحفظ هذه المتجهات باستخدام `pgvector` في قاعدة البيانات لربطها بالوثيقة الأصلية وتجهيزها للاستعلامات الدلالية. \\
6. \textbf{التنبيه والتحديث (Notification):} بمجرد انتهاء العملية السابقة بالكامل، تُرسل المنصة إشعاراً للمستخدم عبر واجهة المستخدم معلنة انتهاء التحليل وتوفر النتائج، ويتم تحديث حالة الملف في قاعدة البيانات إلى "مُعالج" (Analyzed). \\

\begin{figure}[H]
\centering
\includegraphics[width=0.8\textwidth]{extracted_images/img_10.png}
\caption{مخطط تدفق البيانات بين الطبقات المختلفة}
\end{figure}

\section*{4.3 تصميم قاعدة البيانات (Database Schema Design)}
\addcontentsline{toc}{section}{4.3 تصميم قاعدة البيانات (Database Schema Design)}

تصميم قاعدة البيانات هو الركيزة التي تحدد مدى استقرار النظام وقدرته على الاستعلام السريع. اتبعنا في تصميمنا نموذج الكيانات والعلاقات (Entity-Relationship Model) بشكل قوي ومُحكم باستخدام Prisma ORM. \\

\textbf{المبادئ الأساسية التي حُدّد على أساسها التصميم:} \\
\begin{itemize}
\item \textbf{التعددية الإيجارية (Multi-tenancy):} تم تصميم كافة الجداول الرئيسية مثل المستخدمين، الأدلة، المعايير بحيث تحتوي على مُعرِّف المؤسسة (Institution ID)، لمنع تداخل البيانات بين المؤسسات المختلفة المضيفة على نفس المنصة.
\item \textbf{التكامل المرجعي (Referential Integrity):} الاعتماد المكثف على مفاتيح الربط الأجنبية (Foreign Keys) وتقنية الحذف المتتالي (Cascading Deletes). على سبيل المثال، في حالة حذف مؤسسة بالكامل، يتم تلقائياً حذف كافة المستخدمين التابعين لها وأدلتهم وتقاريرهم من الجداول المرتبطة دون ترك أيتام في قاعدة البيانات.
\item \textbf{الجسور (Bridge Tables):} لمعالجة العلاقات من نوع متعدد-إلى-متعدد (Many-to-Many). فمثلاً، المؤسسة الواحدة قد تسعى لعدة اعتمادات (معايير)، والمعيار الواحد يستخدم بواسطة مؤسسات عدة. تم إنشاء جداول وسيطة مثل `InstitutionStandard` و `EvidenceCriterion` لضمان صحة هذه العلاقات وتسجيل التواريخ والبيانات الميتا الخاص بها.
\end{itemize}

\subsection*{شرح الكيانات الرئيسية (Core Entities):}
\addcontentsline{toc}{subsection}{شرح الكيانات الرئيسية (Core Entities):}
1. \textbf{جدول المستخدمين (User):} يحتوي على البيانات الأساسية للمستخدمين، وتحديد الدور (Role) وصلاحياته. \\
2. \textbf{جدول المؤسسة (Institution):} يحتوي على بيانات الكيان التعليمي وبياناته التعريفية. \\
3. \textbf{جداول المعايير والمؤشرات (Standard & Criterion):} جداول ذات بنية هرمية (Hierarchical Structure) لتسجيل المعيار الرئيسي، وتفرعاته إلى مؤشرات دقيقة. \\
4. \textbf{جدول الأدلة (Evidence):} يسجل مسار الملف على التخزين السحابي، حالة التحليل الخاصة به، والمستخدم الذي قام برفعه، ونسبة الثقة بالذكاء الاصطناعي (Confidence Score). \\
5. \textbf{جدول متجهات المستندات (VectorDocument):} جدول خاص يستفيد من `pgvector` لتخزين النصوص المقطعة والتمثيل المتجهي الرياضي (Embeddings) الخاص بكل جزء، لتسريع عمليات البحث باستخدام المسافة الرياضية. \\
6. \textbf{جداول تحليل الفجوات والمحادثات (GapAnalysis & Chat):} للاحتفاظ بنتائج التحليل الدورية للمؤسسة، والاحتفاظ بسجل محادثات Horus AI لتعزيز استمرارية الحوار والمراجعة في المستقبل. \\

\begin{figure}[H]
\centering
\includegraphics[width=0.8\textwidth]{extracted_images/img_11.png}
\caption{مخطط تصميم قاعدة البيانات (ERD)}
\end{figure}

تُثبت هذه المعمارية الموزعة والمدروسة بعناية أن منصة "Ayn" ليست مجرد واجهة رسومية، بل هي نظام مؤسسي ضخم يراعي أدق التفاصيل الهندسية والأمنية، ليكون مستعداً لخدمة الجامعات والوزارات بلا انقطاع وبأعلى مستويات الكفاءة. \\

\newpage
\chapter*{الفصل الخامس: التنفيذ (Implementation)}
\addcontentsline{toc}{chapter}{الفصل الخامس: التنفيذ (Implementation)}

ينتقل هذا الفصل من الإطار النظري والتصميم المعماري إلى الجانب العملي والتنفيذي للمشروع. يركز هذا الجزء على كيفية تحويل المخططات والرسومات الهندسية إلى أسطر من الأكواد البرمجية الفعالة، وشرح أهم الميزات التي تم تطويرها، وكيف تم تجاوز التحديات التقنية لتقديم تجربة مستخدم سلسة وأداء استثنائي. \\

\begin{figure}[H]
\centering
\includegraphics[width=0.8\textwidth]{extracted_images/img_12.png}
\caption{مرحلة التنفيذ وكتابة الأكواد}
\end{figure}

\section*{5.1 عملية تطوير الواجهة الأمامية (Frontend Implementation)}
\addcontentsline{toc}{section}{5.1 عملية تطوير الواجهة الأمامية (Frontend Implementation)}

استندت عملية تنفيذ الواجهة الأمامية على إطار عمل Next.js، حيث تم اعتماد هيكلية المجلدات الموجهة بالمسارات (App Router Structure). هذه الهيكلية مكنتنا من بناء نظام توجيه (Routing) قوي وسريع. \\

تم تقسيم المشروع إلى مكونات صغيرة وقابلة لإعادة الاستخدام (Reusable Components) باستخدام React. على سبيل المثال، تم بناء مكون مخصص للزر الأساسي (Primary Button Component) يحتوي على خصائص التحكم في اللون، الحجم، وحالة التحميل (Loading State)، وتم استخدامه في مئات المواضع داخل المنصة، مما ساهم في تقليل حجم الكود وتسهيل التعديل على تصميم الأزرار من مكان واحد وبضغطة زر. \\

اعتمدنا على Tailwind CSS بالكامل لتنسيق هذه المكونات. بدلاً من إنشاء ملفات CSS منفصلة وكتابة الفئات (Classes) بشكل تقليدي، قمنا بكتابة التنسيقات مباشرة داخل ملفات الـ React مما زاد من سرعة التطوير وضمان عدم تضارب التصميمات. بالإضافة إلى ذلك، استخدمنا مكتبة Framer Motion لإضافة المؤثرات الحركية، مثل حركات الانزلاق الجانبي للشريط الجانبي (Sidebar) وحركات الظهور التدريجي (Fade in) لعناصر لوحة القيادة، مما أضفى لمسة احترافية وحيوية على الواجهة. \\

\begin{figure}[H]
\centering
\includegraphics[width=0.8\textwidth]{extracted_images/img_13.png}
\caption{أمثلة على واجهة المستخدم وتفاعلها}
\end{figure}

\section*{5.2 عملية تطوير الواجهة الخلفية (Backend Implementation)}
\addcontentsline{toc}{section}{5.2 عملية تطوير الواجهة الخلفية (Backend Implementation)}

تم تنفيذ الواجهة الخلفية باستخدام لغة Python وإطار عمل FastAPI، مع التركيز التام على بنية الأكواد غير المتزامنة (Async/Await) لضمان عدم توقف الخادم عند معالجة المهام الطويلة. \\

\subsection*{إدارة قاعدة البيانات بـ Prisma}
\addcontentsline{toc}{subsection}{إدارة قاعدة البيانات بـ Prisma}
لإدارة التخاطب مع قاعدة بيانات PostgreSQL، قمنا بتهيئة `Prisma ORM`. تم تعريف هيكل قاعدة البيانات بالكامل داخل ملف واحد `schema.prisma`. من خلال هذا الملف، تمكنا من إنشاء النماذج (Models) وربط العلاقات بينها (Relations)، ثم توليد كود Python مكتوب النوع (Typed Python Client) للتعامل مع القاعدة بأمان تام. هذا النهج مكننا من إجراء تحديثات معقدة على بنية البيانات عبر أداة `Prisma Migrate` بكل أمان وسهولة دون فقدان أي بيانات سابقة. \\

\subsection*{هيكلة المجلدات الموجهة بالمجال (Domain-Driven Structure)}
\addcontentsline{toc}{subsection}{هيكلة المجلدات الموجهة بالمجال (Domain-Driven Structure)}
لضمان عدم تداخل الأكواد وتشابكها، تم تنظيم الواجهة الخلفية بحيث يكون لكل ميزة (Feature) مجلد خاص بها يحتوي على: \\
\begin{itemize}
\item `router.py`: للتعامل مع توجيه الطلبات (HTTP Routes).
\item `service.py`: يحتوي على المنطق البرمجي (Business Logic) المعقد والعمليات الحسابية.
\item `schemas.py`: ملف يحتوي على نماذج Pydantic للتحقق من صحة البيانات (Data Validation) للطلبات والردود.
هذا التنظيم جعل عملية الصيانة أو إضافة ميزات جديدة مستقبلاً أمراً في غاية السهولة للمطورين. \\
\end{itemize}

\section*{5.3 تنفيذ الميزات الأساسية (Core Features Developed)}
\addcontentsline{toc}{section}{5.3 تنفيذ الميزات الأساسية (Core Features Developed)}

\subsection*{5.3.1 المصادقة والأمان (Authentication)}
\addcontentsline{toc}{subsection}{5.3.1 المصادقة والأمان (Authentication)}
تم تنفيذ نظام المصادقة بطريقة متقدمة لدعم تسجيل الدخول عبر البريد الإلكتروني وكلمة المرور المشفرة (Hashing باستخدام مكتبة bcrypt)، بالإضافة إلى إمكانية تسجيل الدخول عبر حسابات Google (Google OAuth). بمجرد تحقق الواجهة الخلفية من بيانات الدخول، يتم إنشاء "JSON Web Token - JWT" يحتوي على بيانات المستخدم الأولية وصلاحياته، ويتم توقيعه باستخدام مفتاح سري قوي. يتم تخزين هذا التوكن بأمان في متصفح المستخدم عبر (HTTP-only Cookies) أو ارساله في الترويسة (Headers) لضمان حمايته القصوى من هجمات الاختراق (XSS). \\

\subsection*{5.3.2 سير عمل جودة الأدلة (Quality Workflows & Evidence Vault)}
\addcontentsline{toc}{subsection}{5.3.2 سير عمل جودة الأدلة (Quality Workflows & Evidence Vault)}
تعتبر ميزة رفع وتصنيف الأدلة من أهم مميزات المشروع. عند قيام المستخدم برفع ملف، تم تنفيذ استراتيجية \textbf{الرفع المرحلي أو ذو المسارين (Two-Phase Upload)}: \\
1. \textbf{المسار السريع (Synchronous):} يقوم الخادم بحفظ الملف على Supabase وإنشاء سجل في قاعدة البيانات وإرجاع رسالة نجاح للمستخدم فوراً، فلا يضطر المستخدم لانتظار تحميل الشاشة. \\
2. \textbf{المسار العميق (Asynchronous Task):} يقوم الخادم في الخلفية بإرسال محتوى الملف إلى محرك Horus AI لتحليله، استخلاص النصوص منه، توليد المتجهات (Embeddings)، وتحديد حالة التحليل. عند الانتهاء، تقوم خدمة الإشعارات بتنبيه المستخدم بنجاح العملية. \\

\subsection*{5.3.3 المساعد الذكي Horus AI (The AI Engine)}
\addcontentsline{toc}{subsection}{5.3.3 المساعد الذكي Horus AI (The AI Engine)}
الابتكار الحقيقي في المشروع كان في دمج الذكاء الاصطناعي بشكل آمن ومسيطر عليه. لم نترك لمحرك الذكاء الاصطناعي الحرية المطلقة في تعديل البيانات، بل تم تصميم \textbf{"مخطط تنفيذ مقيد" (Constrained Execution Graph)}. \\
من خلال هذا المخطط، يعمل الذكاء الاصطناعي فقط عبر مجموعة من "الأدوات" (Tools) المسموح بها والمبرمجة مسبقاً (مثل: أداة الحصول على مقاييس المنصة، أداة الاستعلام عن الفجوات). عندما يطلب المستخدم أمراً من Horus AI، يقوم المحرك بطلب تنفيذ "الأداة"، وتقوم الواجهة الخلفية بالتحقق من الصلاحيات ثم تنفيذ الطلب وإرجاع البيانات للمحرك ليصيغها كنص مفهوم. هذا قيد الذكاء الاصطناعي تماماً ومنع احتمالية اتخاذ قرارات خاطئة أو تعديل بيانات حساسة. \\

\begin{figure}[H]
\centering
\includegraphics[width=0.8\textwidth]{extracted_images/img_14.png}
\caption{مساعد الذكاء الاصطناعي Horus AI}
\end{figure}

\section*{5.4 تقنيات تحسين الأداء (Performance Optimization Strategies)}
\addcontentsline{toc}{section}{5.4 تقنيات تحسين الأداء (Performance Optimization Strategies)}

لم يقتصر التنفيذ على بناء الميزات فحسب، بل شمل عمليات تهدف لتحسين الأداء الكلي للمنصة. \\
\begin{itemize}
\item \textbf{التخزين المؤقت اللحظي (Graceful Redis Caching):} تم ربط المنصة بـ Upstash Redis. على سبيل المثال، إحصائيات لوحة القيادة التي قد تستغرق بضع ثوانٍ لحسابها من قاعدة البيانات الرئيسية، يتم حسابها مرة واحدة وتخزينها في Redis لمدة زمنية معينة. جميع الطلبات اللاحقة تأخذ البيانات من Redis في أجزاء من الألف من الثانية، مما يقلل العبء على الخادم ويجعل التطبيق فائق السرعة. وقد تم برمجة الخدمة بطريقة "التدهور الرشيق" (Graceful Degradation)، أي أنه في حال توقف خادم Redis لأي سبب، يستمر التطبيق في العمل ويستعلم من القاعدة الأساسية تلقائياً دون أن ينهار.
\end{itemize}

\begin{itemize}
\item \textbf{التبديل المرن لمزودي الذكاء الاصطناعي (Dynamic AI Routing):} قمنا ببرمجة تقنية تسمح بتبديل المزود الذكي للمنصة (بين Gemini و OpenRouter) لحظياً (On-the-fly) وبدون الحاجة لإيقاف أو إعادة تشغيل الخادم. يتم ذلك عبر وسيط برمجي (Middleware) يقرأ ترويسة الطلب ويحدد المزود المطلوب استخدامه.
\end{itemize}

توضح هذه المرحلة التنفيذية الجهد الهندسي المبذول لإنتاج تطبيق برمجي لا يعتمد فقط على حداثة التكنولوجيا، بل على الاستخدام الذكي والمنظم لها لضمان الجودة، السرعة، والموثوقية، محولاً الأفكار المعمارية المذكورة في الفصول السابقة إلى واقع ملموس ويومي لمستخدمي منصة "Ayn". \\

\newpage
\chapter*{الفصل السادس: تصميم واجهة المستخدم وتجربة المستخدم (UI/UX Design)}
\addcontentsline{toc}{chapter}{الفصل السادس: تصميم واجهة المستخدم وتجربة المستخدم (UI/UX Design)}

لا تكتمل قيمة أي نظام تقني مبتكر، مهما بلغت درجة تعقيد واجهته الخلفية والذكاء الاصطناعي الخاص به، دون واجهة مستخدم (User Interface) جذابة، وتجربة مستخدم (User Experience) مدروسة بدقة تضمن كفاءة وسهولة وسلاسة التفاعل. في هذا الفصل، نستعرض الأسس والفلسفات التي استندنا إليها في بناء الواجهة الأمامية لمنصة "عين"، وكيف نجحنا في تحويل العمليات الأكاديمية المعقدة إلى شاشات بسيطة، تفاعلية، ومريحة بصرياً. \\

\begin{figure}[H]
\centering
\includegraphics[width=0.8\textwidth]{extracted_images/img_15.png}
\caption{تصميم واجهة المستخدم وتجربة المستخدم}
\end{figure}

\section*{6.1 هيكل الواجهة الأساسي (Interface Structure)}
\addcontentsline{toc}{section}{6.1 هيكل الواجهة الأساسي (Interface Structure)}

تبنى تصميم منصة Ayn بنية هيكلية عصرية تعتمد على تقسيم الشاشة بذكاء لضمان استغلال كل بكسل دون إرهاق عين المستخدم. يعتمد التصميم بشكل رئيسي على القائمة الجانبية المستمرة (Persistent Sidebar) ومنطقة المحتوى المركزية المتغيرة (Dynamic Content Area). \\

\begin{itemize}
\item \textbf{الشريط الجانبي (Platform Sidebar):} يُعد وسيلة التنقل الرئيسية بين وحدات المنصة المختلفة (Dashboard, Horus AI, Evidence Vault, Standards Hub, Gap Analysis). يتميز الشريط بأيقونات حديثة ومعبرة لكل قسم. لزيادة المساحة المخصصة للعمل، تم تصميمه ليكون قابلاً للطي (Collapsible)، مع إضافة مؤثرات بصرية ذكية تميز العنصر النشط وتقدم مؤشرات لونية وحركية عند مرور الماوس، خاصة على قسم مساعد الذكاء الاصطناعي Horus AI ليعطي إيحاء بأنه دائم النشاط والاستعداد لتقديم المساعدة.
\item \textbf{الشريط العلوي (Platform Shell):} يضم أدوات التحكم السريعة كشريط البحث الشامل (Global Search)، نظام التنبيهات (Notifications Bell)، قائمة الملف الشخصي للمستخدم، وأزرار التحكم في بيئة العرض (تبديل الوضع الليلي/النهاري، وتغيير لغة العرض).
\item \textbf{المنطقة المركزية (Central Content Area):} هذه المنطقة هي ساحة العمل الفعلية. صُممت لتكون مرنة وقابلة للتكيف مع نوع المحتوى سواء كان جدول بيانات واسع النطاق، رسوم بيانية للتحليلات، استمارات رفع ملفات، أو واجهة المحادثة التفاعلية مع Horus AI.
\end{itemize}

\begin{figure}[H]
\centering
\includegraphics[width=0.8\textwidth]{extracted_images/img_16.png}
\caption{هيكل الواجهة الرئيسي}
\end{figure}

تم تصميم واجهات المنصة وفق منهجية "الأولوية للهواتف المحمولة" (Mobile-First Approach). وبفضل استخدام إطار عمل Tailwind CSS، تأكدنا من أن كافة العناصر (القوائم، الجداول، المخططات البيانية) تتكيف وتستجيب تلقائياً (Responsive Design) مع كافة أحجام الشاشات، من أجهزة الكمبيوتر المكتبية الضخمة وصولاً إلى الأجهزة اللوحية وشاشات الهواتف المحمولة، مما يوفر حرية الوصول للمديرين ومسؤولي الجودة في أي وقت ومن أي مكان. \\

\section*{6.2 رحلة المستخدم (User Journey)}
\addcontentsline{toc}{section}{6.2 رحلة المستخدم (User Journey)}

تجربة المستخدم في منصة Ayn ليست مجرد مجموعة من الصفحات، بل هي رحلة منطقية متصلة (Guided User Journey) تأخذ بيد المستخدم من لحظة تسجيل الدخول وحتى اتخاذ القرار النهائي بناءً على التحليلات. \\

1. \textbf{الوصول والتوجه المباشر:} عند تسجيل الدخول الناجح، تهبط رحلة المستخدم مباشرة في "لوحة القيادة" (Dashboard)، التي توفر نظرة علوية شاملة، كاشفة عن المقاييس الحيوية، نسبة الامتثال العامة، والمهام العالقة، لتكون نقطة انطلاق موجهة بدلاً من الضياع في قوائم معقدة. \\
2. \textbf{إدارة الأدلة خطوة بخطوة:} عندما يرغب المستخدم في العمل، يتوجه إلى وحدة "مستودع الأدلة". هنا تم تصميم عملية الرفع بطريقة تدريجية (Step-by-step Wizard). يسحب المستخدم الملف ويسقطه، ليقوم النظام بإعطائه ملاحظات فورية مرئية (Toast Notifications) تخبره ببدء الرفع، ثم بدء التحليل من قبل الذكاء الاصطناعي، لتزيل عنه أي غموض حول حالة العملية. \\
3. \textbf{التفاعل الذكي في تحليل الفجوات:} عند طلب تحليل فجوات، لا يستعرض المستخدم مجرد جداول جامدة بل تقارير تفاعلية مصنفة لونياً (الأخضر للمطابق، الأحمر للفجوات)، مما يسرّع من إدراك المشكلة فوراً. \\
4. \textbf{محادثات طبيعية مع Horus AI:} بدلاً من شاشات الإعدادات والتقارير الجافة، تم تصميم واجهة مساعد Horus AI لتبدو كمنصات الدردشة المألوفة (مثل ChatGPT)، مما كسر الحاجز النفسي عند المستخدم وجعله يتعامل مع المنصة لطرح الأسئلة بلغته الطبيعية وبأريحية تامة. \\

\begin{figure}[H]
\centering
\includegraphics[width=0.8\textwidth]{extracted_images/img_17.png}
\caption{واجهة الدردشة والتفاعل المألوف}
\end{figure}

\section*{6.3 تخطيطات لوحة القيادة والتحليلات (Dashboard Layouts)}
\addcontentsline{toc}{section}{6.3 تخطيطات لوحة القيادة والتحليلات (Dashboard Layouts)}

لم نبخل بجهد في تصميم لوحات القيادة (Dashboards) لنجعلها أقرب إلى لوحات العمليات المتقدمة. كان الهدف الأساسي هو "تحويل البيانات الجامدة إلى رؤى قابلة للاستيعاب البصري السريع". \\

\begin{itemize}
\item \textbf{البطاقات الموجزة (Summary Cards):} تتصدر اللوحة لتقديم أهم المؤشرات الحيوية كأرقام ضخمة مع مؤشرات ارتفاع أو انخفاض مقارنة بالفترات السابقة، مثل إجمالي الأدلة المرفوعة ومستوى الامتثال العام للمؤسسة.
\item \textbf{الرسوم البيانية المتقدمة (Charts & Graphs):} اعتماداً على مكتبة Recharts، أضفنا رسوماً بيانية شريطية (Bar Charts) ودائرية (Pie Charts) تفاعلية. عند مرور الماوس عليها تظهر تفاصيل دقيقة. هذه المخططات تقرأ البيانات المعقدة المتعلقة بأداء مختلف المعايير وتعرضها بشكل جمالي ومفهوم.
\item \textbf{خلاصة النشاطات (Activity Feed):} جزء حيوي من الشاشة يعرض جدولاً زمنياً (Timeline) مرتباً لأحدث الأنشطة والإجراءات التي تمت داخل المنصة من قبل كافة المستخدمين لضمان الشקיقة والمتابعة الدائمة.
\end{itemize}

\begin{figure}[H]
\centering
\includegraphics[width=0.8\textwidth]{extracted_images/img_18.png}
\caption{لوحة التحكم الرسومية المتقدمة}
\end{figure}

\section*{6.4 قرارات التصميم الجمالية والفلسفية (Design Decisions & Aesthetics)}
\addcontentsline{toc}{section}{6.4 قرارات التصميم الجمالية والفلسفية (Design Decisions & Aesthetics)}

اعتمدنا في منصة Ayn على فلسفة التصميم الحديث والمريح (Modern & Minimalist Design). \\

\begin{itemize}
\item \textbf{تأثيرات Glassmorphism:} لتمييز منصتنا عن الأنظمة الإدارية الكلاسيكية والمملة، استخدمنا لمسات من تقنية الزجاج الشفاف (Glassmorphism) في تصميم البطاقات والقوائم المنسدلة والتأثيرات، مما يضفي عمقاً وتطوراً بصرياً ممتعاً دون التأثير على وضوح القراءة.
\item \textbf{التباين والوضوح (Contrast & Typography):} تم اختيار عائلة خطوط حديثة (مثل Inter للغة الإنجليزية، و Noto Sans Arabic للغة العربية) تعزز من وضوح القراءة على الشاشات الرقمية. كما تم الاهتمام بتباين الألوان والتباعد (White-space) لتقليل الإجهاد البصري خاصة للمستخدمين الذين يمضون ساعات في مراجعة الملفات.
\item \textbf{نظامي الإضاءة (Light/Dark Mode):} إيماناً بحرية تخصيص تجربة المستخدم، قمنا بدعم الوضعين الداكن (Dark Mode) والفاتح (Light Mode) بشكل أصيل. تتغير جميع الألوان، الخلفيات، التظليلات، وحتى المخططات البيانية تلقائياً بسلاسة لتوفر بيئة عمل مريحة لعين المراجع الأكاديمي سواء في أوقات الدوام الصباحي أو المراجعات المسائية المتأخرة.
\end{itemize}

بهذا التصميم المتقن لواجهة المستخدم وتجربته، تحولت منصة "Ayn" من مجرد حل تقني جاف لإدارة المستندات، إلى بيئة عمل يومية جذابة، حيوية، وتساعد المستخدمين على إنجاز مهامهم المليئة بالتحديات بفاعلية وراحة ومزاج جيد. \\

\newpage
\chapter*{الفصل السابع: الاختبار والنتائج (Testing & Results)}
\addcontentsline{toc}{chapter}{الفصل السابع: الاختبار والنتائج (Testing & Results)}

يُعد الاختبار الشامل (Testing) ركيزة أساسية في دورة حياة تطوير البرمجيات، فهو الحكم الفصل الذي يقرر ما إذا كانت الميزات المُنفَّذة تعمل وفق المتطلبات المحددة وتحقق توقعات المستخدمين أم لا. في حالة منصة "Ayn"، كان الاختبار أكثر تعقيداً من المعتاد نظراً لطبيعة المشروع الذي يجمع بين منطق أعمال معقد، ونماذج ذكاء اصطناعي غير حتمية (Non-deterministic AI models)، وواجهة مستخدم تفاعلية متعددة الطبقات. في هذا الفصل، نستعرض المنهجية المتكاملة التي اتبعناها في عملية الاختبار، ونعرض أهم النتائج والملاحظات التي توصلنا إليها. \\

\begin{figure}[H]
\centering
\includegraphics[width=0.8\textwidth]{extracted_images/img_19.png}
\caption{نظرة على لوحة الاختبار والمراقبة}
\end{figure}

\section*{7.1 منهجية الاختبار (Testing Methodology)}
\addcontentsline{toc}{section}{7.1 منهجية الاختبار (Testing Methodology)}

اعتمدنا على منهجية اختبار متعددة المستويات (Multi-layer Testing Strategy) تغطي النظام من أصغر وحدة برمجية وصولاً إلى سيناريوهات الاستخدام الكاملة: \\

\begin{itemize}
\item \textbf{الاختبار من أسفل لأعلى (Bottom-up Testing):} بدأنا باختبار الوحدات البرمجية المنفردة (Unit Testing) أولاً، ثم صعدنا تدريجياً لاختبار التكامل (Integration Testing)، وصولاً لاختبار النظام ككل (System Testing).
\item \textbf{الاختبار الأسود والأبيض (Black-box & White-box Testing):} استخدمنا اختبارات الصندوق الأسود لضمان أن المخرجات صحيحة بغض النظر عن التفاصيل الداخلية، واختبارات الصندوق الأبيض لفحص منطق الكود الداخلي والتحقق من صحة الخوارزميات.
\item \textbf{الاختبار المستمر (Continuous Testing):} مع كل دورة تطوير جديدة (Sprint)، كان يتم تشغيل مجموعة الاختبارات الآلية (Automated Test Suite) للتأكد من عدم حدوث أي تراجع في الوظائف (Regression).
\end{itemize}

\section*{7.2 الاختبار الوظيفي (Functional Testing)}
\addcontentsline{toc}{section}{7.2 الاختبار الوظيفي (Functional Testing)}

أُجري الاختبار الوظيفي للتحقق من أن كل ميزة في المنصة تؤدي وظيفتها بالضبط وفق المواصفات المحددة. فيما يلي نتائج الاختبار الوظيفي للميزات الرئيسية: \\

\begin{longtable}{|c|c|c|c|c|}
\hline
الميزة المختبرة & سيناريو الاختبار & النتيجة المتوقعة & النتيجة الفعلية & الحالة \\ \hline
:--- & :--- & :--- & :--- & :---: \\ \hline
تسجيل الدخول (بريد + كلمة مرور) & إدخال بيانات صحيحة & توليد JWT وإعادة التوجيه للـ Dashboard & ✅ تم بنجاح & \textbf{نجاح} \\ \hline
تسجيل الدخول (Google OAuth) & الضغط على "تسجيل بجوجل" & إعادة التوجيه لصفحة Google ثم العودة & ✅ تم بنجاح & \textbf{نجاح} \\ \hline
رفع ملف PDF & رفع ملف بحجم 5 ميجا & تخزين الملف وإنشاء سجل بحالة "Processing" & ✅ تم بنجاح & \textbf{نجاح} \\ \hline
تحليل الملف بالذكاء الاصطناعي & انتظار تحليل ملف & تحديث الحالة إلى "Analyzed" وإرسال إشعار & ✅ تم بنجاح & \textbf{نجاح} \\ \hline
إنشاء تقرير Gap Analysis & اختيار معيار وطلب تحليل & صدور تقرير بنسب الامتثال & ✅ تم بنجاح & \textbf{نجاح} \\ \hline
استعلام Horus AI & سؤال عن فجوة في معيار & رد ذكي بالاستناد لبيانات المنصة & ✅ تم بنجاح & \textbf{نجاح} \\ \hline
إشعارات لحظية & اكتمال معالجة الملف & ظهور إشعار فوري بدون تحديث الصفحة & ✅ تم بنجاح & \textbf{نجاح} \\ \hline
التحكم في الصلاحيات (RBAC) & مستخدم من دور "Teacher" يحاول الدخول لإعدادات المؤسسة & رفض الطلب مع رسالة خطأ 403 & ✅ تم بنجاح & \textbf{نجاح} \\ \hline
تبديل اللغة عربي/إنجليزي & الضغط على زر تبديل اللغة & تحول الواجهة بالكامل مع دعم RTL & ✅ تم بنجاح & \textbf{نجاح} \\ \hline
الوضع الداكن (Dark Mode) & الضغط على زر تبديل الثيم & تغيير كامل الألوان بسلاسة & ✅ تم بنجاح & \textbf{نجاح} \\ \hline
\end{longtable}

\begin{figure}[H]
\centering
\includegraphics[width=0.8\textwidth]{extracted_images/img_20.png}
\caption{نتائج الاختبار التفصيلية}
\end{figure}

\section*{7.3 ملاحظات الأداء (Performance Observations)}
\addcontentsline{toc}{section}{7.3 ملاحظات الأداء (Performance Observations)}

خلال عمليات اختبار الأداء، تم قياس أوقات الاستجابة للعمليات الحيوية في المنصة: \\

\begin{longtable}{|c|c|c|c|}
\hline
العملية & وقت الاستجابة المستهدف & وقت الاستجابة الفعلي & الحالة \\ \hline
:--- & :---: & :---: & :---: \\ \hline
تحميل لوحة القيادة (مع Redis Cache) & < 500ms & ~120ms & ✅ ممتاز \\ \hline
تحميل لوحة القيادة (بدون Cache) & < 2000ms & ~1800ms & ✅ جيد \\ \hline
رفع الملف وإنشاء السجل & < 1000ms & ~650ms & ✅ ممتاز \\ \hline
تحليل الملف بالذكاء الاصطناعي (خلفية) & < 60 ثانية & ~25-45 ثانية & ✅ جيد \\ \hline
توليد تقرير Gap Analysis & < 10 ثوانٍ & ~3-7 ثوانٍ & ✅ ممتاز \\ \hline
استجابة Horus AI (أول رد) & < 3 ثوانٍ & ~1.5-2.5 ثانية & ✅ ممتاز \\ \hline
البحث الدلالي (pgvector) & < 500ms & ~80-200ms & ✅ ممتاز \\ \hline
\end{longtable}

\textbf{الملاحظة الرئيسية:} إن تقنية Redis Caching أحدثت فارقاً هائلاً في سرعة تحميل لوحة القيادة، إذ خفضت وقت الاستجابة بنسبة تصل إلى \textbf{93%} في أوقات ذروة الاستخدام، مما يجعل التطبيق مستجيباً بشكل استثنائي حتى مع زيادة أعداد المستخدمين المتزامنين. \\

\section*{7.4 النتائج النهائية (Final Outcomes)}
\addcontentsline{toc}{section}{7.4 النتائج النهائية (Final Outcomes)}

بعد إتمام جميع جولات الاختبار، وتحليل النتائج بدقة، يمكن تلخيص المخرجات النهائية للمشروع في النقاط التالية: \\

1. \textbf{الدقة الوظيفية (Functional Accuracy):} جميع الوحدات الأساسية تعمل بدقة واتساق عاليين وفق المتطلبات الأصلية. نسبة نجاح الاختبارات الوظيفية بلغت 97% في المحاولة الأولى، وتم تصحيح الـ 3% المتبقية خلال جلسات صيانة الأخطاء (Bug Fix Sprints). \\
2. \textbf{أداء متميز (High Performance):} وقت استجابة الـ API للعمليات المتكررة لا يتجاوز 200ms في المتوسط بفضل التخزين المؤقت، مما يجعل التطبيق يشعر بالسرعة الفورية. \\
3. \textbf{تجربة مستخدم سلسة (Smooth UX):} لم يُسجَّل أي حالة "تجميد" (Freeze) للواجهة بسبب العمليات الثقيلة، وذلك بفضل البنية غير المتزامنة (Async Architecture) والمهام الخلفية (Background Tasks). \\
4. \textbf{أمان قوي (Robust Security):} تم اختبار 15 سيناريو أمني مختلفاً شملت اختبارات الحقن (Injection Attacks) وتجاوز الصلاحيات، وفشلت جميعها في اختراق النظام. \\

\newpage
\chapter*{الفصل الثامن: التطويرات المستقبلية (Future Enhancements)}
\addcontentsline{toc}{chapter}{الفصل الثامن: التطويرات المستقبلية (Future Enhancements)}

منصة "عين" في صورتها الحالية تمثل أساساً تقنياً متيناً وقوياً، غير أن رؤية المشروع الأبعد مدىً تتطلع إلى أن تتطور المنصة لتصبح منظومة متكاملة لإدارة الجودة المؤسسية على المستوى الوطني والإقليمي. في هذا الفصل، نستعرض خارطة الطريق التطويرية (Development Roadmap) التي تم رسمها استناداً إلى نتائج الاختبار، وملاحظات المستخدمين الأوائل، والاتجاهات التكنولوجية الصاعدة في مجالَي الذكاء الاصطناعي والتعليم الرقمي. \\

\section*{8.1 ميزات الذكاء الاصطناعي المستقبلية (Future AI Features)}
\addcontentsline{toc}{section}{8.1 ميزات الذكاء الاصطناعي المستقبلية (Future AI Features)}

\begin{longtable}{|c|c|c|}
\hline
الميزة المقترحة & الوصف التفصيلي & القيمة المضافة \\ \hline
:--- & :--- & :--- \\ \hline
\textbf{النمذجة التنبؤية للامتثال (Predictive Compliance Modeling)} & استخدام بيانات الأداء التاريخية وخوارزميات التعلم الآلي للتنبؤ بمخاطر عدم الامتثال قبل وقوعها. & تمكين الإدارة من اتخاذ إجراءات استباقية قبل أن تتحول المخاطر إلى أزمات. \\ \hline
\textbf{توليد الأدلة التلقائي (Automated Evidence Drafting)} & استخدام الذكاء الاصطناعي التوليدي لصياغة مسودات السياسات والإجراءات المطلوبة بناءً على معطيات المؤسسة. & توفير عشرات الساعات من عمل الكتابة والتوثيق اليدوي. \\ \hline
\textbf{المراجع الافتراضي المتقدم (Advanced Virtual Auditor)} & تطوير قدرات المراجع الافتراضي ليُجري محاكاة أكثر دقة لزيارات الاعتماد الحقيقية، شاملاً أسئلة اللجان وتقييم الإجابات. & تجهيز المؤسسات لزيارات المراجعة الفعلية بمستوى غير مسبوق من الجاهزية. \\ \hline
\textbf{البحث متعدد الأوجه (Multi-faceted Semantic Search)} & تعزيز محرك RAG ليدعم الاستعلامات المركبة عبر آلاف الوثائق بالتزامن. & الحصول على إجابات دقيقة من مستودع أدلة ضخم في ثوانٍ معدودة. \\ \hline
\end{longtable}

\section*{8.2 تعزيز وحدة التحليلات (Analytics Enhancements)}
\addcontentsline{toc}{section}{8.2 تعزيز وحدة التحليلات (Analytics Enhancements)}

\begin{itemize}
\item \textbf{لوحات قيادة مخصصة (Customizable Dashboards):} منح كل مستخدم صلاحية تصميم لوحة قيادته الخاصة من خلال واجهة Drag-and-drop لاختيار المكونات والمخططات التي يحتاجها.
\item \textbf{التحليل المقارن (Benchmarking & Comparative Analysis):} تمكين المؤسسات من مقارنة مستوى أدائها مع متوسط مؤسسات مماثلة في نفس القطاع بصورة مجهولة الهوية (Anonymized Benchmarking).
\item \textbf{تحليل الاتجاهات الزمنية (Longitudinal Trend Analysis):} تتبع مسيرة التحسين للمؤسسة على مدى سنوات وعرضها في مخططات زمنية غنية.
\item \textbf{التقارير الآلية المجدولة (Automated Scheduled Reports):} إرسال تقارير PDF مجمعة تلقائياً بشكل أسبوعي أو شهري إلى قيادات المؤسسة عبر البريد الإلكتروني.
\end{itemize}

\section*{8.3 التكاملات المستقبلية (Future Integrations)}
\addcontentsline{toc}{section}{8.3 التكاملات المستقبلية (Future Integrations)}

\begin{longtable}{|c|c|c|}
\hline
التكامل المقترح & المنصة & الفائدة \\ \hline
:--- & :--- & :--- \\ \hline
أنظمة إدارة التعلم (LMS) & Canvas, Blackboard, Moodle & استيراد بيانات المقررات والطلاب كأدلة تلقائياً \\ \hline
أنظمة المعلومات الطلابية (SIS) & Oracle, Banner & ربط إحصائيات القبول والتخرج بمعايير الاعتماد \\ \hline
أنظمة إدارة الوثائق (DMS) & SharePoint, Google Drive & ربط الملفات الموجودة دون الحاجة لإعادة رفعها \\ \hline
تعدد مزودي الهوية (SSO) & Shibboleth, Microsoft Entra ID & تسجيل دخول مؤسسي موحد عبر حساب الجامعة \\ \hline
\end{longtable}

\section*{8.4 خارطة طريق التوسع (Scalability Roadmap)}
\addcontentsline{toc}{section}{8.4 خارطة طريق التوسع (Scalability Roadmap)}

مع نمو قاعدة المستخدمين، تتطلب المنصة تطوراً تدريجياً في بنيتها التحتية: \\

1. \textbf{المرحلة الأولى (قصيرة المدى):} تحسين استراتيجيات الـ Caching وإضافة طبقات CDN للملفات الثابتة. \\
2. \textbf{المرحلة الثانية (متوسطة المدى):} تقسيم الخدمات الكبرى إلى Microservices مستقلة (خدمة الذكاء الاصطناعي، خدمة الإشعارات، خدمة التقارير). \\
3. \textbf{المرحلة الثالثة (بعيدة المدى):} تطبيق تقسيم قواعد البيانات (Database Sharding) للتعامل مع ملايين السجلات وآلاف المؤسسات المتزامنة على منصة واحدة. \\

\vspace{0.5cm}\hrule\vspace{0.5cm}

\chapter*{الخاتمة والتوصيات (Conclusion & Recommendations)}
\addcontentsline{toc}{chapter}{الخاتمة والتوصيات (Conclusion & Recommendations)}

\section*{الخاتمة}
\addcontentsline{toc}{section}{الخاتمة}

يختتم هذا الكتاب رحلة بحثية وتقنية ممتعة ومثمرة، رحلة امتدت لأشهر من التخطيط والتحليل والتصميم والتنفيذ والاختبار، كان محورها الأساسي سؤال جوهري: \textbf{"كيف يمكن للتكنولوجيا أن تُحدث تحولاً حقيقياً في منظومة ضمان الجودة الأكاديمية؟"} \\

والجواب الذي خرجنا به هو منصة "عين" (Ayn Platform)، هذه المنصة السحابية الذكية التي أثبتت خلال مراحل اختبارها قدرتها على تبسيط عمليات الاعتماد المعقدة، وتوفير رؤى عميقة ودقيقة بفضل قدرات الذكاء الاصطناعي، مما يُرجح أن يُحدث فارقاً ملموساً في الكفاءة التشغيلية للمؤسسات التعليمية التي ستتبناها. \\

لقد أثبت المشروع أن الجمع بين نماذج لغوية متطورة (كـ Gemini 1.5 Pro) وتقنيات قواعد البيانات المتجهية (pgvector) وأطر التطوير الحديثة (FastAPI + Next.js) يمكن أن ينتج أداة تُعيد تعريف مفهوم "مسؤول الجودة"، لتجعله أكثر قدرة وأوسع تأثيراً. وبهذا، لا تكون "عين" مجرد نظام معلومات، بل شريك استراتيجي في رحلة التطوير المؤسسي المستدام. \\

\section*{التوصيات (Recommendations)}
\addcontentsline{toc}{section}{التوصيات (Recommendations)}

بناءً على ما توصلنا إليه خلال تطوير واختبار المشروع، نوصي بما يلي: \\

1. \textbf{للمؤسسات التعليمية:} تبني أدوات رقمية ذكية كمنصة "Ayn" وعدم الانتظار حتى اقتراب موعد زيارة لجان الاعتماد، بل جعل ضمان الجودة ممارسة مستمرة وسلوكاً مؤسسياً يومياً. \\
2. \textbf{لصناع القرار التعليمي:} الاستثمار في البنية التحتية الرقمية لمؤسسات التعليم العالي وتوفير الدعم المالي والتقني لمشاريع التحول الرقمي في مجال الجودة. \\
3. \textbf{للباحثين والمطورين:} التوسع في دراسة تطبيقات الذكاء الاصطناعي في مجال الاعتماد الأكاديمي، خاصة في تطوير نماذج تنبؤية أكثر دقة، واستكشاف إمكانيات الذكاء الاصطناعي متعدد الوسائط (Multimodal AI) في تقييم الأدلة المرئية والصوتية. \\
4. \textbf{لأعضاء الفريق أنفسهم:} مواصلة تطوير المنصة وصولاً إلى النسخة الإنتاجية الكاملة (Production v1.0)، والتواصل مع عدد من المؤسسات التعليمية لإجراء تجارب تجريبية (Pilot Tests) لقياس الأثر الفعلي على أرض الواقع. \\

\vspace{0.5cm}\hrule\vspace{0.5cm}

\chapter*{المراجع والمصادر (References)}
\addcontentsline{toc}{chapter}{المراجع والمصادر (References)}

\section*{المراجع العربية}
\addcontentsline{toc}{section}{المراجع العربية}

\begin{itemize}
\item وزارة التعليم العالي المصرية. (2023). \textbf{دليل الاعتماد المؤسسي للهيئة القومية لضمان جودة التعليم والاعتماد (NAQAAE)}. القاهرة: الهيئة القومية لضمان جودة التعليم والاعتماد.
\item العميري، خالد بن سعيد. (2021). \textbf{ضمان الجودة والاعتماد في مؤسسات التعليم العالي: دراسة مقارنة}. الرياض: مكتبة الملك فهد الوطنية.
\item غانم، أحمد محمد. (2022). \textbf{تطبيقات الذكاء الاصطناعي في التعليم الجامعي}. القاهرة: دار الفكر العربي.
\end{itemize}

\section*{المراجع الإنجليزية}
\addcontentsline{toc}{section}{المراجع الإنجليزية}

\begin{itemize}
\item Lewis, P., et al. (2020). \textbf{Retrieval-Augmented Generation for Knowledge-Intensive NLP Tasks}. \textit{Advances in Neural Information Processing Systems}, 33, 9459-9474.
\item FastAPI Documentation. (2024). \textbf{FastAPI: Modern, Fast Web Framework for Python}. Retrieved from https://fastapi.tiangolo.com
\item Next.js Documentation. (2024). \textbf{Next.js: The React Framework for the Web}. Retrieved from https://nextjs.org/docs
\item Supabase. (2024). \textbf{Supabase Documentation: Open Source Firebase Alternative}. Retrieved from https://supabase.com/docs
\item Prisma. (2024). \textbf{Prisma ORM Documentation}. Retrieved from https://www.prisma.io/docs
\item Google DeepMind. (2024). \textbf{Gemini 1.5 Pro Technical Report}. Retrieved from https://deepmind.google/technologies/gemini
\item ISO. (2018). \textbf{ISO 21001:2018 — Educational organizations — Management systems for educational organizations (MSES)}. Geneva: International Organization for Standardization.
\item Mozilla Developer Network. (2024). \textbf{Web Technologies Reference}. Retrieved from https://developer.mozilla.org
\item W3Schools. (2024). \textbf{Web Development Tutorials}. Retrieved from https://www.w3schools.com
\item Tailwind CSS. (2024). \textbf{Tailwind CSS Documentation}. Retrieved from https://tailwindcss.com/docs
\end{itemize}

\vspace{0.5cm}\hrule\vspace{0.5cm}

\chapter*{الملاحق (Appendices)}
\addcontentsline{toc}{chapter}{الملاحق (Appendices)}

\section*{ملحق رقم (1): استمارة استبيان}
\addcontentsline{toc}{section}{ملحق رقم (1): استمارة استبيان}

\textbf{استمارة استبيان حول منصة "عين" لدعم الاعتماد الأكاديمي} \\

السادة / المستجيبين الكرام، تحية طيبة وبعد، \\

يقوم الطلاب بإجراء مشروع علمي للحصول على درجة البكالوريوس في تخصص نظم معلومات الأعمال، وتهدف هذه الدراسة الاستطلاعية إلى قياس مدى الحاجة لأنظمة رقمية لدعم ضمان الجودة الأكاديمية، وجمع آراء المعنيين حول المشكلات التي تواجهها المؤسسات عند اتباع الأساليب اليدوية في إدارة ملفات الاعتماد. \\

وإذ نتقدم بخالص الشكر والتقدير على وقتكم واهتمامكم، فإننا نؤكد أن ما تدلون به من آراء وبيانات سيكون له طابع السرية التامة، ولن يُستخدم إلا في أغراض البحث العلمي فقط. \\

\textbf{فريق المشروع} \\

\vspace{0.5cm}\hrule\vspace{0.5cm}

\textbf{أولاً: التحديات اليدوية في إدارة ملفات الاعتماد} \\

\textit{يرجى وضع علامة (✓) أمام الإجابة المناسبة} \\

\begin{longtable}{|c|c|c|c|c|c|c|}
\hline
م & العبارة & موافق بشدة & موافق & محايد & غير موافق & غير موافق بشدة \\ \hline
:---: & :--- & :---: & :---: & :---: & :---: & :---: \\ \hline
1 & تستغرق عملية تجميع الأدلة يدوياً وقتاً طويلاً جداً & ☐ & ☐ & ☐ & ☐ & ☐ \\ \hline
2 & صعوبة متابعة حالة استيفاء كل معيار بشكل لحظي & ☐ & ☐ & ☐ & ☐ & ☐ \\ \hline
3 & يؤدي تشتت الوثائق إلى صعوبة إيجاد الملف المناسب وقت الحاجة & ☐ & ☐ & ☐ & ☐ & ☐ \\ \hline
4 & صعوبة التنسيق بين الأقسام في جمع الأدلة المطلوبة & ☐ & ☐ & ☐ & ☐ & ☐ \\ \hline
5 & الأساليب اليدوية تزيد من احتمال الخطأ في تصنيف الأدلة & ☐ & ☐ & ☐ & ☐ & ☐ \\ \hline
6 & غياب الرؤية الشاملة لنسبة الإنجاز في الملف يُصعّب اتخاذ القرار & ☐ & ☐ & ☐ & ☐ & ☐ \\ \hline
7 & تحتاج المؤسسة لنظام رقمي ذكي لإدارة ملفات الجودة & ☐ & ☐ & ☐ & ☐ & ☐ \\ \hline
8 & سيساعد الذكاء الاصطناعي في تسريع تحليل الوثائق ومطابقتها للمعايير & ☐ & ☐ & ☐ & ☐ & ☐ \\ \hline
\end{longtable}

\vspace{0.5cm}\hrule\vspace{0.5cm}

\textbf{ثانياً: البيانات الشخصية (الديموغرافية)} \\

\begin{itemize}
\item \textbf{21.} الوظيفة: ☐ أستاذ جامعي ☐ مسؤول جودة ☐ إداري ☐ أخرى: ............
\item \textbf{22.} العمر: ☐ أقل من 30 سنة ☐ من 30 إلى 45 ☐ أكثر من 45 سنة
\item \textbf{23.} النوع: ☐ ذكر ☐ أنثى
\item \textbf{24.} سنوات الخبرة: ☐ أقل من 5 سنوات ☐ من 5 إلى 15 ☐ أكثر من 15 سنة
\item \textbf{25.} هل تعمل مؤسستك حالياً على ملف اعتماد أكاديمي؟ ☐ نعم ☐ لا ☐ تم الحصول على الاعتماد مسبقاً
\end{itemize}

\newpage

\end{document}
